\documentclass[11pt]{article}

\usepackage[utf8]{inputenc}
% microtype 의 폰트 확장은 scalable 폰트를 요구한다. 기본 CM 이 비트맵(Type3)으로
% 잡히는 배포판에서는 lmodern 없이는 "auto expansion is only possible with
% scalable fonts" 로 컴파일이 죽는다.
\usepackage{lmodern}
\usepackage[T1]{fontenc}
\usepackage[margin=1in]{geometry}
\usepackage{amsmath,amssymb}
\usepackage{graphicx}
\usepackage{longtable,booktabs}
\usepackage{microtype}
% svgnames 를 줘야 teal 이 정의된다. 없으면 \cite 하나마다
% "Undefined color `teal'" 이 뜬다 (본편 산출물에서 732회).
\usepackage[svgnames]{xcolor}
% hyperref goes last; it makes \cite clickable through to the bibliography.
\usepackage[colorlinks=true,linkcolor=blue,citecolor=teal,urlcolor=magenta]{hyperref}

% pandoc emits \tightlist but expects the preamble to define it.
\providecommand{\tightlist}{%
  \setlength{\itemsep}{0pt}\setlength{\parskip}{0pt}}

% If the compile times out on Overleaf, comment out \tableofcontents below first.
\title{Comprehensive Survey on Instruction Tuning for Large Language Models}
\author{Generated by AutoSurvey}
\date{}

\begin{document}
\maketitle
\tableofcontents
\newpage

\section{Introduction to Instruction Tuning}

\subsection{Definition and Significance of Instruction Tuning}

Instruction tuning is a crucial technique used to enhance the performance of large language models (LLMs) by aligning them with user intentions. At its core, instruction tuning involves fine-tuning pre-trained LLMs on a dataset of instruction-response pairs, where the instructions are designed to elicit specific responses from the model \cite{ref1}. This process allows the model to learn the patterns and relationships between instructions and responses, ultimately enabling it to generalize to new, unseen instructions. By doing so, instruction tuning enables LLMs to understand and follow instructions provided by users, thereby improving their ability to generate accurate and contextually appropriate responses.

The significance of instruction tuning lies in its ability to unlock the full potential of LLMs, enabling them to perform a wide range of tasks, from simple text generation to complex problem-solving \cite{ref2}. Moreover, instruction tuning has been shown to improve the robustness and reliability of LLMs, making them more suitable for real-world applications where accuracy and consistency are critical \cite{ref3}. One of the key benefits of instruction tuning is its ability to adapt LLMs to specific tasks and domains, allowing for the creation of highly specialized and effective models in areas such as translation, summarization, or question-answering \cite{ref4}.

In addition to its practical applications, instruction tuning also has significant implications for our understanding of LLMs and their capabilities. By studying how LLMs respond to instructions and learn to follow them, researchers can gain insights into the underlying mechanisms and representations that drive these models \cite{ref5}. This can help to inform the development of new models and techniques, ultimately leading to more advanced and capable LLMs \cite{ref6}. Furthermore, instruction tuning can be used to create LLMs that are specifically designed for tasks such as conversation, dialogue systems, or natural language understanding \cite{ref7}.

Despite its many benefits, instruction tuning is not without its challenges and limitations. One of the primary challenges is the need for high-quality instruction data, which can be time-consuming and expensive to create \cite{ref8}. Moreover, instruction tuning can be sensitive to the specific instructions and responses used in the fine-tuning dataset, which can affect the model's performance and robustness \cite{ref3}. To address these challenges, researchers have developed various techniques and strategies for creating and selecting instruction data, such as data filtering and augmentation \cite{ref9}.

As the field of instruction tuning continues to evolve, there is a growing interest in exploring its potential for multilingual and multimodal applications. For example, researchers have developed techniques for instruction tuning that can be applied to multiple languages \cite{ref10} and modalities \cite{ref11}. These advances have significant implications for the development of more advanced and versatile LLMs that can be applied to a wide range of tasks and domains. In the next section, we will delve into the history and current state of instruction tuning, highlighting key milestones, breakthroughs, and challenges that have shaped the field into its current form.

\subsection{Brief History and Current State of Instruction Tuning}

The concept of instruction tuning has undergone significant evolution since its inception, with key milestones marking its progression from fine-tuning pre-trained language models on specific tasks to enabling large language models to follow human instructions and adapt to new tasks \cite{ref12}. A crucial breakthrough in this journey was the introduction of the concept of ``instruction-following'' abilities in large language models \cite{ref13}, which paved the way for the development of various methods and techniques for instruction tuning, including supervised fine-tuning, reinforcement learning, and prompt engineering.

As the field advanced, researchers began to explore more sophisticated instruction tuning methods, such as adaptive task balancing and intra-task difficulty assessment \cite{ref14}, to improve the efficiency and effectiveness of the process. Recent advancements have focused on optimizing instruction synthesis \cite{ref15} and selecting the most informative and relevant instruction data \cite{ref16}. Furthermore, there is a growing interest in exploring the potential of instruction tuning in multimodal and multilingual settings \cite{ref17}, as well as its applications in real-world tasks such as code generation, dialogue systems, and natural language understanding \cite{ref18}.

The current state of instruction tuning is characterized by a growing recognition of its importance in enabling large language models to generalize to unseen tasks and adapt to new instructions \cite{ref19}. To achieve this, researchers are actively exploring new methods and techniques, including the use of transfer learning, meta-learning, and few-shot learning. Moreover, there is an increasing focus on evaluating the performance of instruction-tuned models on a wide range of tasks and benchmarks \cite{ref20}. Despite the significant progress made, challenges such as the quality and diversity of instruction data \cite{ref21} and the need for more robust and generalizable instruction tuning methods \cite{ref22} remain to be addressed.

In the context of the broader field of large language models, instruction tuning has emerged as a critical technique for unlocking their full potential. By enabling these models to understand and follow instructions, instruction tuning has opened up new avenues for applications in areas such as translation, summarization, and question-answering \cite{ref4}. As researchers continue to push the boundaries of what is possible with instruction tuning, we can expect to see significant improvements in the performance and generalizability of large language models \cite{ref23}. The ongoing advancements in instruction tuning are expected to play a crucial role in shaping the future of large language models and their applications in various domains.

\subsection{Scope and Organization of the Survey}

This survey on instruction tuning for large language models aims to provide a comprehensive overview of the current state of research in this field, building on the foundation established by the history and evolution of instruction tuning. The scope of this survey is broad, covering various aspects of instruction tuning, including its definition, significance, and applications, as well as the methods and techniques used to enable large language models to follow human instructions and adapt to new tasks \cite{ref12}. The survey is organized into several sections, each focusing on a specific topic related to instruction tuning, and is designed to inform and guide researchers and practitioners in their efforts to develop more effective and efficient instruction tuning methods.

The survey begins with an introduction to instruction tuning, including its definition, significance, and brief history, which provides a foundation for understanding the concept of instruction tuning and its importance in the context of large language models. The introduction is followed by a section on background and fundamentals, which covers the basics of large language models, their architecture, training methods, and the principles of instruction tuning, including the concept of ``instruction-following'' abilities in large language models \cite{ref13}.

The survey then delves into the methods and techniques used for instruction tuning, including supervised fine-tuning, reinforcement learning, and prompt engineering, which have been developed to improve the performance of large language models on specific tasks \cite{ref24}. This section provides an in-depth analysis of the various approaches used to tune large language models, highlighting their strengths and weaknesses, as well as recent advancements in instruction tuning, such as adaptive task balancing and intra-task difficulty assessment \cite{ref14}. The applications and evaluations of instruction tuning are also discussed, including its use in natural language understanding, generation, and multimodal tasks, and the potential of instruction tuning in real-world applications, such as code generation, dialogue systems, and natural language understanding \cite{ref18}.

The survey also addresses the challenges and limitations of instruction tuning, such as data quality, catastrophic forgetting, and ethical concerns, which are critical to the development of more general, efficient, and safe large language models \cite{ref22}. This section provides a critical analysis of the obstacles that researchers and practitioners face when implementing instruction tuning, highlighting the need for further research and development, and the importance of addressing these challenges to enable large language models to generalize to unseen tasks and adapt to new instructions \cite{ref19}. A comparison of instruction tuning methods is also presented, categorizing and analyzing different approaches based on their performance, efficiency, and applicability, and discussing the potential applications of instruction tuning in emerging areas such as multimodal learning and edge AI.

The survey concludes with a discussion on future directions and open research questions in instruction tuning, including the potential applications of instruction tuning in various domains, and the implications of instruction tuning for the development of more general, efficient, and safe large language models \cite{ref23}. The final section provides recommendations for practitioners and researchers, outlining potential avenues for future research, and highlighting the importance of instruction tuning in enabling large language models to follow human instructions and adapt to new tasks. Throughout the survey, we highlight key findings and contributions from existing research, including the importance of instruction tuning in improving the performance of large language models, the need for more efficient and effective methods for instruction tuning, and the potential applications of instruction tuning in various domains.

\section{Background and Fundamentals}

\subsection{Introduction to Large Language Models}

Large language models (LLMs) have revolutionized the field of natural language processing (NLP) in recent years, enabling significant advancements in various tasks and applications. At the core of LLMs is the concept of training a single model on a massive dataset of text, allowing it to learn the patterns and structures of language \cite{ref8}. This approach has led to the development of powerful models that can process and understand human language, generating coherent and contextually relevant text.

The architecture of LLMs is typically based on the transformer model, which was introduced in \cite{ref11}. The transformer model uses self-attention mechanisms to weigh the importance of different words in a sentence, allowing the model to capture long-range dependencies and contextual relationships. This architecture has been widely adopted in LLMs, including popular models such as BERT, RoBERTa, and XLNet \cite{ref25}. The transformer model's ability to capture complex patterns and relationships in language has made it an essential component of LLMs.

LLMs are typically trained using a combination of supervised and unsupervised learning methods. The supervised learning approach involves training the model on a labeled dataset, where the model is given a input text and a corresponding output text. The unsupervised learning approach involves training the model on a large corpus of text, where the model is given a input text and asked to predict the next word in the sequence \cite{ref2}. This approach allows the model to learn the patterns and structures of language, without the need for labeled data. As a result, LLMs can be fine-tuned for specific tasks, enabling them to achieve state-of-the-art results in various NLP tasks.

One of the key challenges in training LLMs is the need for large amounts of computational resources and data. Training a single LLM can require thousands of hours of computation time and millions of dollars of investment \cite{ref26}. To address this challenge, researchers have developed various techniques, including distributed training, where the model is trained on multiple machines in parallel, and transfer learning, where a pre-trained model is fine-tuned on a smaller dataset \cite{ref27}. These techniques have enabled the development of more efficient and effective LLMs, which can be applied to a wide range of tasks and applications.

Despite the challenges, LLMs have achieved state-of-the-art results in various NLP tasks, including language translation, text summarization, and question answering \cite{ref28}. These models have also been used in various applications, including chatbots, language translation software, and text generation tools \cite{ref7}. However, LLMs also have some limitations, including the need for large amounts of data and computational resources, and the risk of overfitting to the training data \cite{ref29}. To address these limitations, researchers have developed various techniques, including data augmentation, where the training data is augmented with additional examples, and regularization techniques, where the model is penalized for overfitting to the training data \cite{ref30}.

The development of LLMs has also led to a growing interest in instruction tuning, which involves fine-tuning a pre-trained LLM on a specific task or dataset \cite{ref8}. Instruction tuning has been shown to improve the performance of LLMs on various tasks, including language translation, text summarization, and question answering \cite{ref31}. This approach has also been used in various applications, including chatbots, language translation software, and text generation tools \cite{ref7}. As the field of NLP continues to evolve, instruction tuning is likely to play an increasingly important role in the development of more accurate and effective LLMs. The next section will delve into the fundamentals of instruction tuning, including key concepts, techniques, and methodologies \cite{ref32}.

\subsection{Fundamentals of Instruction Tuning}

Instruction tuning is a crucial technique used to enhance the performance of large language models (LLMs) by fine-tuning them on instruction-response pairs. As discussed in the previous section, LLMs have revolutionized the field of natural language processing, and instruction tuning has emerged as a key approach to improve their performance on specific tasks. The primary goal of instruction tuning is to align the model with human intentions and preferences, enabling it to generate more accurate and contextually appropriate responses.

At its core, instruction tuning involves training a pre-trained LLM on a dataset of instruction-response pairs, where each instruction is a natural language prompt that specifies a task or question, and the corresponding response is the desired output. This process enables the model to learn the patterns and relationships between instructions and responses, allowing it to generalize to new, unseen instructions \cite{ref24}. By doing so, instruction tuning can significantly improve the model's ability to understand and respond to user requests, making it a vital component of many NLP applications.

One of the key concepts in instruction tuning is the idea of ``instructability,'' which refers to the model's ability to follow instructions and generate accurate responses \cite{ref5}. Instructability is critical in real-world applications, where models are expected to perform tasks that are specified by users through natural language instructions. To achieve high instructability, models must be able to understand the nuances of language, including context, intent, and semantics \cite{ref2}. This requires a deep understanding of the relationships between language, context, and task-specific requirements.

Several techniques and methodologies have been proposed to improve the effectiveness of instruction tuning. One popular approach is to use a combination of supervised and reinforcement learning, where the model is trained on instruction-response pairs and then fine-tuned using reinforcement learning to optimize its performance on specific tasks \cite{ref31}. This approach can help to improve the model's ability to generalize to new tasks and instructions, while also reducing the need for large amounts of labeled training data.

In addition to these techniques, researchers have also explored the use of different evaluation metrics and benchmarks to assess the performance of instruction-tuned models. For example, the ``Sensitivity'' metric has been proposed to evaluate the model's sensitivity to variations in instructions \cite{ref33}. These evaluation metrics can help to identify the strengths and weaknesses of different instruction tuning approaches, providing valuable insights for future research and development.

Despite the progress made in instruction tuning, there are still several challenges and limitations that need to be addressed. One of the major challenges is the need for high-quality instruction-response pairs, which can be time-consuming and expensive to obtain \cite{ref21}. To address these challenges, researchers have proposed several methodologies, including data selection and filtering techniques, to improve the quality and diversity of the instruction-response pairs \cite{ref34}. These methodologies can help to reduce the cost and effort required to develop high-quality instruction-response pairs, making it more feasible to apply instruction tuning to a wide range of NLP tasks and applications.

In conclusion, instruction tuning is a powerful technique for enhancing the performance of large language models. By fine-tuning models on instruction-response pairs, researchers can improve the model's instructability, enabling it to generate more accurate and contextually appropriate responses. As we will discuss in the next section, the architecture and training methods for instruction tuning have been a subject of extensive research in recent years, with various approaches being proposed to improve the effectiveness and efficiency of instruction tuning \cite{ref32}.

\subsection{Architecture and Training Methods for Instruction Tuning}

The architecture and training methods for instruction tuning have been a subject of extensive research in recent years. As large language models (LLMs) continue to advance, the need for effective instruction tuning methods has become increasingly important. In this subsection, we will delve into the various architecture and training methods specifically designed for instruction tuning.

One of the key challenges in instruction tuning is the need to balance the trade-off between model performance and computational resources. To address this challenge, researchers have proposed various architecture designs, such as the Mixture-of-Experts (MoE) model \cite{ref25}. The MoE model is a type of neural network architecture that consists of multiple expert models, each of which is responsible for a specific task or domain. By using a MoE model, researchers can selectively activate or deactivate different expert models during instruction tuning, thereby reducing the computational resources required.

Another important aspect of instruction tuning is the training method used. Traditional training methods, such as supervised fine-tuning, can be effective but often require large amounts of labeled data. To address this limitation, researchers have proposed alternative training methods, such as reinforcement learning from human feedback \cite{ref24}. This approach involves training the model using a reward signal that is based on human feedback, rather than a fixed label. By using reinforcement learning, researchers can train models that are more robust and adaptable to different tasks and domains.

In addition to these architecture and training methods, researchers have also explored the use of data-efficient instruction tuning methods. One such approach is the use of submodular data mixture strategies \cite{ref35}, which involve selecting a subset of the most informative data points to use for instruction tuning. This approach can significantly reduce the amount of data required for instruction tuning, while still achieving comparable performance to traditional methods.

The use of instruction tuning has also been explored in the context of federated learning \cite{ref36}. Federated learning involves training models on decentralized data, where each client has its own private data. By using instruction tuning, researchers can train models that are more robust and adaptable to different tasks and domains, while also preserving the privacy of the clients' data.

Furthermore, instruction tuning can be applied to various applications, including natural language processing, computer vision, and multimodal learning. For instance, instruction tuning can be used to improve the performance of large language models in language translation \cite{ref2}, text summarization \cite{ref37}, and question answering \cite{ref38}. Additionally, instruction tuning can be used in low-resource settings \cite{ref39} and in the context of continual learning \cite{ref40}.

In conclusion, instruction tuning is a powerful technique that can be used to improve the performance of large language models in many tasks and applications. By leveraging various architecture and training methods, researchers can develop more effective and efficient instruction tuning approaches that can be applied to a wide range of domains and applications. As the field continues to evolve, we can expect to see even more innovative and effective approaches to emerge, which can further improve the performance of large language models in many areas.

\section{Methods and Techniques for Instruction Tuning}

\subsection{Supervised Fine-Tuning and Reinforcement Learning}

Supervised fine-tuning and reinforcement learning from human feedback are two fundamental methods employed in the instruction tuning of large language models (LLMs). To understand the significance of these methods, it is essential to recognize the importance of fine-tuning in adapting LLMs to specific tasks and datasets. Fine-tuning enables LLMs to learn from labeled data or human preferences, allowing them to generate more accurate and informative responses. Supervised fine-tuning involves training a model on a specific task with labeled data, where the model learns to predict the correct output based on the input. This approach has been widely used in various natural language processing tasks, including text classification, sentiment analysis, and machine translation. On the other hand, reinforcement learning from human feedback is a technique that utilizes human preferences to fine-tune LLMs. This method involves training a reward model to predict the human preference between two model outputs and then using this reward model to fine-tune the LLM.

The emergence of large language models (LLMs) \cite{ref41} has led to significant advancements in natural language processing. However, these models often require fine-tuning to align with human preferences and values. Supervised fine-tuning is a common approach used to fine-tune LLMs, but it has its limitations. For instance, supervised fine-tuning requires large amounts of labeled data, which can be time-consuming and expensive to obtain. Moreover, supervised fine-tuning may not always capture the nuances of human preferences, leading to suboptimal performance. In contrast, reinforcement learning from human feedback \cite{ref42} has been proposed as an alternative to supervised fine-tuning. This approach involves training a reward model to predict human preferences and then using this reward model to fine-tune the LLM.

One of the key advantages of reinforcement learning from human feedback is its ability to capture nuanced human preferences. Unlike supervised fine-tuning, which relies on labeled data, reinforcement learning from human feedback can learn from human preferences and adapt to changing preferences over time. Moreover, reinforcement learning from human feedback can be used to fine-tune LLMs on a wide range of tasks, including text generation, dialogue systems, and language translation. For example, \cite{ref43} demonstrates the use of reinforcement learning from human feedback to fine-tune a language model to be more helpful and harmless. Despite its advantages, reinforcement learning from human feedback also has its challenges. For instance, \cite{ref44} highlights the memory consumption issues associated with reinforcement learning from human feedback. Moreover, \cite{ref45} discusses the challenges of using reinforcement learning from human feedback in fine-tuning language models.

In recent years, there has been a growing interest in combining supervised fine-tuning and reinforcement learning from human feedback. For example, \cite{ref46} proposes a supervised iterative learning approach that combines supervised fine-tuning and reinforcement learning from human feedback. This approach has been shown to be effective in fine-tuning LLMs and achieving state-of-the-art results on various natural language processing tasks. In addition to combining supervised fine-tuning and reinforcement learning from human feedback, researchers have also explored other techniques to improve the instruction tuning of LLMs. For instance, \cite{ref47} proposes a fine-grained human feedback approach that provides more detailed and nuanced feedback to the model. This approach has been shown to be effective in improving the performance of LLMs on various tasks. Another technique that has been proposed is the use of natural language feedback \cite{ref48}. This approach involves providing feedback to the model in the form of natural language, rather than relying on labeled data or human preferences.

As the field of natural language processing continues to evolve, it is likely that we will see further advancements in the instruction tuning of LLMs. The development of more efficient and effective fine-tuning methods, such as those discussed in the following section, will play a crucial role in improving the performance of LLMs. By combining supervised fine-tuning and reinforcement learning from human feedback, as well as exploring other techniques such as fine-grained human feedback and natural language feedback, researchers can create more accurate and informative LLMs that are better aligned with human values and preferences \cite{ref49}. Ultimately, the goal of instruction tuning is to create LLMs that can generate high-quality responses that are both helpful and harmless, and ongoing research in this area is bringing us closer to achieving this goal.

\subsection{Fine-Tuning Methods and Parameter-Efficient Techniques}

Fine-tuning methods are crucial for instruction tuning, as they enable large language models to adapt to specific tasks and datasets. In this subsection, we will discuss various fine-tuning methods used for instruction tuning, including full fine-tuning, parameter-efficient fine-tuning, and transfer learning. These methods are essential for achieving excellent performance on target tasks, and their selection depends on the specific task, dataset, and computational resources available.

Full fine-tuning involves updating all the parameters of a pre-trained model on a specific task or dataset. This approach can lead to excellent performance on the target task but is often computationally expensive and requires a large amount of training data \cite{ref50}. Moreover, full fine-tuning can result in catastrophic forgetting, where the model forgets the knowledge it learned during pre-training \cite{ref51}. To mitigate these issues, researchers have proposed various techniques, such as regularization and knowledge distillation, which can help alleviate catastrophic forgetting \cite{ref52}.

To address the limitations of full fine-tuning, parameter-efficient fine-tuning methods have been proposed. These methods update only a small subset of the model's parameters, reducing the computational cost and memory requirements \cite{ref53}. One popular approach is adapter-based fine-tuning, which involves adding small, trainable modules to the pre-trained model \cite{ref54}. Another approach is prefix-tuning, which involves adding a small set of learnable parameters to the model's input embeddings \cite{ref55}. These methods have been shown to be effective in achieving excellent performance on target tasks while minimizing computational costs and memory requirements.

Transfer learning is also a widely used fine-tuning method, which involves pre-training a model on a large dataset and then fine-tuning it on a smaller, task-specific dataset \cite{ref56}. This approach can be particularly effective when the target task has limited training data. However, transfer learning can also suffer from catastrophic forgetting, and techniques such as regularization and knowledge distillation can be used to mitigate this issue \cite{ref52}. Recently, there has been a growing interest in developing more efficient and effective fine-tuning methods, such as sparse fine-tuning methods, which involve updating only a subset of the model's parameters \cite{ref57}.

In addition to these methods, various techniques can be used to improve the efficiency and effectiveness of fine-tuning. For example, gradual tuning approaches involve fine-tuning the model's parameters in a progressive manner \cite{ref58}. Another technique is to use knowledge distillation, which involves transferring knowledge from a pre-trained model to a smaller, task-specific model \cite{ref59}. Moreover, \cite{ref60} proposed a method for generating a sparse mask in a task-agnostic manner, which can be used to fine-tune a model on multiple tasks without introducing new latency. This approach can be particularly useful in scenarios where models need to be deployed on edge devices or in real-time applications.

In conclusion, fine-tuning methods are a crucial component of instruction tuning, and various approaches have been proposed to improve their efficiency and effectiveness. By selecting the most suitable fine-tuning method and technique, it is possible to achieve excellent performance on a wide range of tasks while minimizing computational costs and memory requirements. The choice of fine-tuning method depends on the specific task, dataset, and computational resources available, and future research should focus on developing more efficient and effective fine-tuning methods, as well as exploring new techniques for improving the performance of instruction-tuned models \cite{ref61}. The development of more efficient and effective fine-tuning methods will have a significant impact on the field of natural language processing, enabling the creation of more accurate and informative large language models that can be deployed in a wide range of applications.

\subsection{Prompt Engineering and Optimization Techniques}

Prompt engineering and optimization techniques are essential for improving instruction tuning for large language models, as they enable the design and optimization of prompts to elicit specific responses from the model, thereby enhancing its performance on various tasks. Building on the fine-tuning methods discussed earlier, prompt engineering and optimization techniques provide a complementary approach to adapting pre-trained language models to downstream tasks. In this subsection, we will explore the different prompt engineering and optimization techniques used to improve instruction tuning, and examine how they can be used in conjunction with fine-tuning methods to achieve optimal results.

One of the key techniques used in prompt engineering is prompt tuning, which involves optimizing the prompt embeddings to adapt the pre-trained language model to downstream tasks \cite{ref62}. This approach has been shown to be effective in improving the performance of large language models on various natural language processing tasks. Another technique is prompt optimization, which involves optimizing the prompt to maximize the model's performance on a specific task \cite{ref63}. By combining prompt tuning and optimization, researchers can develop more effective instruction tuning methods that can improve the performance of large language models on a wide range of tasks.

Recent studies have also explored the use of reinforcement learning to optimize prompts for instruction tuning \cite{ref64}. This approach involves using a reinforcement learning algorithm to optimize the prompt based on the model's performance on a specific task. Other techniques, such as gradient-based optimization and evolutionary algorithms, have also been used to optimize prompts for instruction tuning \cite{ref65}. These techniques can be used to optimize prompts for specific tasks and models, and can be combined with other prompt engineering and optimization techniques to achieve even better results.

In addition to these techniques, researchers have also explored the use of multimodal prompts, which involve combining text and image prompts to improve the model's performance on multimodal tasks \cite{ref66}. This approach has been shown to be effective in improving the performance of large language models on tasks such as visual question answering and image captioning. Furthermore, the design of effective prompt templates is also an important aspect of prompt engineering, as a good prompt template should be able to elicit a specific response from the model, while also being flexible enough to accommodate different input formats and tasks \cite{ref67}.

The evaluation of prompt engineering and optimization techniques is also a crucial aspect of instruction tuning, as it enables researchers to compare the performance of different techniques and identify the most effective ones for specific tasks and models. Researchers have proposed various evaluation metrics, such as perplexity and accuracy, to measure the effectiveness of different prompt engineering and optimization techniques \cite{ref21}. By using these metrics, researchers can develop more effective instruction tuning methods that can improve the performance of large language models on a wide range of tasks.

In conclusion, prompt engineering and optimization techniques play a vital role in improving instruction tuning for large language models, and can be used in conjunction with fine-tuning methods to achieve optimal results. By exploring different prompt engineering and optimization techniques, such as prompt tuning, prompt optimization, and multimodal prompts, researchers can develop more effective instruction tuning methods that can improve the performance of large language models on a wide range of tasks \cite{ref68}. Furthermore, the use of reinforcement learning and evolutionary algorithms to optimize prompts for instruction tuning has shown promising results \cite{ref69}, and the development of effective prompt engineering and optimization techniques is an active area of research, with new techniques and methods being proposed and evaluated continuously \cite{ref70}.

\section{Applications and Evaluations of Instruction Tuning}

\subsection{Natural Language Understanding and Generation Tasks}

Instruction tuning has emerged as a crucial technique for enhancing the performance of large language models (LLMs) on various natural language understanding and generation tasks. By aligning the model's outputs with human preferences, instruction tuning enables LLMs to comprehend and respond to human instructions effectively. This technique has been shown to improve the performance of LLMs on various natural language understanding tasks, such as question answering, sentiment analysis, and text classification \cite{ref71,ref2}. For instance, the work presented in \cite{ref71} demonstrates that instruction tuning can significantly enhance the zero-shot performance of LLMs on unseen tasks, including natural language understanding tasks.

In addition to natural language understanding tasks, instruction tuning has also been applied to natural language generation tasks, such as text summarization, dialogue generation, and language translation \cite{ref72,ref73}. The work presented in \cite{ref72} introduces a novel framework for generating knowledge-intensive multi-turn dialogues for instruction tuning, which enables LLMs to produce more accurate and contextually nuanced responses. Furthermore, the study presented in \cite{ref73} proposes a two-stage instruction tuning framework for evaluating the text in both seen and unseen aspects, which demonstrates the effectiveness of instruction tuning in improving the performance of LLMs on natural language generation tasks.

The application of instruction tuning extends beyond natural language understanding and generation tasks to multimodal tasks, which involve processing and generating multiple forms of data, such as text, images, and audio \cite{ref66,ref74}. The work presented in \cite{ref66} introduces a novel approach to multimodal prompt tuning, which enables LLMs to effectively integrate visual and textual prompts during fine-tuning. Moreover, the study presented in \cite{ref74} proposes a dataset comprising 2.8 million multimodal instruction-response pairs, which demonstrates the effectiveness of instruction tuning in improving the performance of LLMs on multimodal tasks.

Instruction tuning has also been applied to various downstream tasks, such as code generation, sentiment analysis, and question answering \cite{ref75,ref76}. The work presented in \cite{ref75} introduces a domain-specific instruction dataset for biomedical natural language processing, which demonstrates the effectiveness of instruction tuning in improving the performance of LLMs on biomedical tasks. Furthermore, the study presented in \cite{ref76} proposes a novel instruction tuning dataset for information retrieval tasks, which demonstrates the effectiveness of instruction tuning in improving the performance of LLMs on search tasks.

The effectiveness of instruction tuning in improving the performance of LLMs on various tasks can be attributed to its ability to align the model's outputs with human preferences \cite{ref24,ref38}. The work presented in \cite{ref24} introduces a simple yet effective task selection method for instruction tuning, which demonstrates the importance of selecting relevant tasks for improving the performance of LLMs. Moreover, the study presented in \cite{ref38} proposes a novel in-context instruction tuning method, which demonstrates the effectiveness of using examples to improve the performance of LLMs on instruction tuning tasks.

In conclusion, instruction tuning has been widely applied to various natural language understanding and generation tasks, including multimodal tasks, and has demonstrated significant improvements in the performance of LLMs on these tasks \cite{ref77,ref78}. The work presented in \cite{ref77} introduces a novel concept of compositional instructions, which enables LLMs to solve complex instructions composed of multiple subtasks. Furthermore, the study presented in \cite{ref78} proposes a novel framework for dialogue NLU, which demonstrates the effectiveness of instruction tuning in improving the performance of LLMs on dialogue tasks. As the field of instruction tuning continues to evolve, it is essential to evaluate the performance of instruction-tuned models using various evaluation metrics and benchmarks, which will be discussed in the following section.

\subsection{Evaluation Metrics and Benchmarks for Instruction Tuning}

Evaluating the performance of instruction-tuned models is a crucial step in understanding their capabilities and limitations. As instruction tuning has been widely applied to various natural language understanding and generation tasks, including multimodal tasks, it is essential to assess the performance of these models using various evaluation metrics and benchmarks. The previous section discussed the applications of instruction tuning, including its ability to improve the performance of large language models on natural language understanding and generation tasks, as well as its potential in multimodal tasks.

One of the key challenges in evaluating instruction-tuned models is the lack of a standardized evaluation metric. Different studies have proposed various metrics, such as perplexity, accuracy, and F1-score, to evaluate the performance of these models \cite{ref79}. However, these metrics may not be suitable for evaluating the performance of instruction-tuned models, as they may not capture the nuances of instruction following. To address this challenge, several benchmarks have been proposed to evaluate the performance of instruction-tuned models, including the PromptSource benchmark \cite{ref12} and the Alpaca benchmark \cite{ref80}.

In addition to these benchmarks, human evaluation metrics, such as human preference and fluency, can provide a more nuanced understanding of the performance of instruction-tuned models \cite{ref81}. These metrics can help to identify the strengths and weaknesses of instruction-tuned models and provide insights into their performance on specific tasks. Furthermore, several studies have proposed new evaluation metrics and methods for evaluating instruction-tuned models, including the Sharpe score \cite{ref82} and the VisIT-Bench benchmark \cite{ref83}.

The FollowEval benchmark \cite{ref84} is another example of a comprehensive evaluation metric, which assesses the performance of instruction-tuned models on five critical dimensions of instruction following, including string manipulation, commonsense reasoning, logical reasoning, spatial reasoning, and response constraints. This benchmark can provide a more comprehensive understanding of the performance of instruction-tuned models and their ability to follow instructions. Other studies, such as \cite{ref85} and \cite{ref86}, have also proposed evaluation metrics and methods for assessing the performance of instruction-tuned models.

The evaluation of instruction-tuned models is a complex task that requires a comprehensive understanding of the strengths and weaknesses of these models. As the field of instruction tuning continues to evolve, it is essential to develop more effective evaluation metrics and benchmarks to assess the performance of these models. The study \cite{ref87} highlights the importance of considering temporal understanding when evaluating instruction-tuned models, and proposes a new benchmark, called TIME, which evaluates the performance of instruction-tuned models on temporal-sensitive video understanding tasks.

In conclusion, evaluating the performance of instruction-tuned models is a crucial step in understanding their capabilities and limitations. Various evaluation metrics and benchmarks have been proposed to assess the performance of these models, including perplexity, accuracy, F1-score, human preference, and fluency. New evaluation metrics and methods, such as the Sharpe score and the VisIT-Bench benchmark, have also been proposed to provide a more nuanced understanding of the performance of instruction-tuned models. These evaluation metrics and benchmarks can help to identify the strengths and weaknesses of instruction-tuned models and provide insights into their performance on specific tasks, which will be further explored in the following section on task-specific instruction tuning \cite{ref16}.

\subsection{Task-Specific Instruction Tuning for Improved Performance}

Task-specific instruction tuning has emerged as a crucial approach to enhance the performance of large language models (LLMs) on specific tasks, building on the foundation of evaluation metrics and benchmarks discussed earlier. This method involves fine-tuning LLMs on task-specific instruction datasets, which enables them to develop a deeper understanding of the task requirements and improve their overall performance. By doing so, task-specific instruction tuning addresses some of the challenges associated with evaluating instruction-tuned models, such as the lack of standardized evaluation metrics and the need for more nuanced understanding of model performance.

One of the primary advantages of task-specific instruction tuning is its ability to adapt LLMs to specific domains or tasks. For instance, \cite{ref24} demonstrates that strategically selecting and training on related tasks can enhance efficiency and prevent performance degradation from learning irrelevant tasks. This approach is particularly useful when dealing with tasks that require specialized knowledge or domain-specific expertise. By fine-tuning LLMs on task-specific instruction datasets, developers can create models that are tailored to specific applications, such as language translation, question answering, or text summarization.

Task-specific instruction tuning can also be used to improve the performance of LLMs on tasks that require complex reasoning or problem-solving skills. For example, \cite{ref88} shows that fine-tuning LLMs on a diverse set of tasks can improve their ability to generalize to new tasks and develop more robust problem-solving skills. This approach is particularly useful for tasks that require creative thinking, critical thinking, or analytical reasoning. By exposing LLMs to a wide range of tasks and instruction datasets, developers can create models that are more versatile and better equipped to handle complex tasks.

Another benefit of task-specific instruction tuning is its ability to reduce the risk of overfitting or underfitting. When LLMs are fine-tuned on large, generic datasets, they may become overly specialized to the training data and fail to generalize to new tasks or domains. Task-specific instruction tuning helps to mitigate this risk by providing LLMs with a more focused and relevant set of training data. For instance, \cite{ref16} proposes a method for selecting high-quality instruction data that is tailored to specific tasks, which can help to improve the performance of LLMs and reduce the risk of overfitting.

Despite its benefits, task-specific instruction tuning also poses several challenges. One of the primary challenges is the need for high-quality, task-specific instruction datasets. Creating these datasets can be time-consuming and labor-intensive, requiring significant expertise and resources. Additionally, task-specific instruction tuning may require significant computational resources, particularly when dealing with large LLMs or complex tasks. \cite{ref89} proposes a method for efficient fine-tuning of LLMs using multiple data sources, which can help to reduce the computational costs associated with task-specific instruction tuning.

To address these challenges, researchers have proposed several methods for improving the efficiency and effectiveness of task-specific instruction tuning. For example, \cite{ref35} introduces a submodular data mixture strategy that can help to optimize the selection of task-specific instruction datasets. This approach can help to reduce the computational costs associated with task-specific instruction tuning while improving the performance of LLMs. Another approach is to use transfer learning or meta-learning techniques to adapt LLMs to new tasks or domains. For instance, \cite{ref90} proposes a method for scaling language model instruction meta-learning, which can help to improve the performance of LLMs on new tasks or domains.

In conclusion, task-specific instruction tuning is a powerful approach for improving the performance of LLMs on specific tasks, and its applications will be further explored in the following subsection. By fine-tuning LLMs on task-specific instruction datasets, developers can create models that are tailored to specific applications, improve their ability to generalize to new tasks, and reduce the risk of overfitting or underfitting. While task-specific instruction tuning poses several challenges, researchers have proposed several methods for improving its efficiency and effectiveness, including the use of submodular data mixture strategies, transfer learning, and meta-learning techniques. As the field of natural language processing continues to evolve, task-specific instruction tuning is likely to play an increasingly important role in the development of more accurate, efficient, and effective LLMs. \cite{ref91} demonstrates that even a small amount of high-quality instruction data can be sufficient for achieving good performance, which highlights the potential of task-specific instruction tuning for improving the efficiency and effectiveness of LLMs.

\section{Challenges and Limitations of Instruction Tuning}

\subsection{Data Quality and Instruction Data Synthesis}

Data quality plays a vital role in instruction tuning, as it directly impacts the performance and reliability of large language models (LLMs). The need for high-quality instruction data is evident, as low-quality data can lead to poor model performance, increased training time, and even catastrophic forgetting \cite{ref13}. To address this issue, it is essential to understand the challenges related to data quality and develop effective methods for instruction data synthesis, selection, and filtering.

One of the primary challenges in instruction tuning is the scarcity of high-quality instruction data. Many existing datasets are limited in size, scope, or quality, which can hinder the performance of LLMs \cite{ref8}. To mitigate this issue, researchers have proposed various methods for instruction data synthesis, including the use of generative models \cite{ref92} and data augmentation techniques \cite{ref30}. These methods aim to generate high-quality instruction data that can supplement existing datasets and improve model performance.

In addition to data scarcity, another challenge related to data quality is the issue of data redundancy. Many instruction datasets contain redundant or similar examples, which can lead to overfitting and decreased model performance \cite{ref93}. To address this issue, researchers have proposed methods for data selection and filtering, such as the use of clustering algorithms \cite{ref94} and diversity-based sampling \cite{ref95}. These methods aim to select a representative subset of high-quality instruction data that can improve model performance while reducing training time and computational resources.

The quality of instruction data is also closely related to the concept of instruction complexity. Instruction complexity refers to the level of difficulty or nuance required to understand and follow an instruction \cite{ref24}. High-quality instruction data should be able to capture a range of instruction complexities, from simple to complex, to enable LLMs to learn and generalize effectively. Researchers have proposed methods for instruction complexity assessment, such as the use of cognitive models \cite{ref21} and machine learning algorithms \cite{ref37}.

Furthermore, the presence of noise or errors in instruction data can also significantly impact model performance. Noisy or erroneous instruction data can lead to poor model performance, as LLMs may learn to replicate errors or biases present in the training data \cite{ref96}. To address this issue, researchers have proposed methods for data cleaning and preprocessing, such as the use of data validation algorithms \cite{ref97} and human evaluation \cite{ref98}.

In the context of federated learning and distributed training paradigms, high-quality instruction data is crucial for ensuring effective knowledge sharing and model performance \cite{ref99}. Researchers have proposed methods for federated instruction tuning, such as the use of decentralized data selection algorithms \cite{ref99} and personalized model training \cite{ref100}. By prioritizing data quality and leveraging advances in machine learning and natural language processing, researchers can improve the performance and reliability of LLMs and enable their widespread adoption in real-world applications \cite{ref32}.

In conclusion, data quality is a critical aspect of instruction tuning, and high-quality instruction data is essential for achieving reliable and effective LLM performance. By addressing the challenges related to data quality, including scarcity, redundancy, complexity, and noise, researchers can develop effective methods for instruction data synthesis, selection, and filtering, ultimately leading to improved model performance and reliability. As instruction tuning continues to evolve, the importance of high-quality instruction data will only continue to grow, making it essential to prioritize data quality in the development of LLMs.

\subsection{Catastrophic Forgetting and Knowledge Retention}

Catastrophic forgetting is a significant challenge in instruction tuning, where models forget previously learned tasks or knowledge when fine-tuned on new tasks. This phenomenon occurs when the model's parameters are updated to fit the new task, causing it to lose its ability to perform well on previous tasks. \cite{ref101} demonstrates that catastrophic forgetting is a common issue in large language models, even when using parameter-efficient fine-tuning methods.

One of the primary reasons for catastrophic forgetting is the interference between old and new knowledge. When a model is fine-tuned on a new task, the new knowledge may overwrite the old knowledge, causing the model to forget its previous capabilities. To mitigate this issue, researchers have proposed various methods, including the use of instruction vectors to capture the model's representations related to specific instruction-following capabilities \cite{ref102}. This approach allows the model to retain its original knowledge while adapting to new tasks.

Another approach to addressing catastrophic forgetting is to use regularization techniques to prevent the model from overwriting its old knowledge. \cite{ref103} proposes a switching mechanism to route computations to parameter-efficient tuned models, reducing interference between old and new knowledge. Additionally, data augmentation and replay can also be used to prevent catastrophic forgetting. \cite{ref104} demonstrates that replaying old data can help mitigate catastrophic forgetting, while \cite{ref105} proposes a method to generate synthetic data that can help mitigate catastrophic forgetting.

Modifying the fine-tuning process itself can also help address catastrophic forgetting. \cite{ref106} proposes a half fine-tuning approach, where half of the model's parameters are frozen to retain old knowledge, while the other half are updated to adapt to new tasks. Knowledge distillation is another technique that can be used to prevent catastrophic forgetting. \cite{ref107} proposes a method to distill knowledge from old tasks to new tasks, reducing the impact of catastrophic forgetting.

The problem of catastrophic forgetting is not unique to instruction tuning and has been explored in other areas of machine learning, such as class-incremental learning. \cite{ref108} proposes a method to remember previous class representations to mitigate catastrophic forgetting in class-incremental learning. Similarly, \cite{ref109} proposes a method to use prompts to mitigate catastrophic forgetting in continual learning.

In the context of instruction tuning, catastrophic forgetting can have significant consequences, such as reducing the model's ability to perform well on previous tasks. To address this issue, researchers have proposed various methods, including \cite{ref110}, \cite{ref111}, and \cite{ref112}. These methods demonstrate the effectiveness of addressing catastrophic forgetting in instruction tuning.

Furthermore, understanding the causes of catastrophic forgetting is crucial to developing effective mitigation strategies. \cite{ref113} proposes a method to anatomy of catastrophic forgetting, demonstrating the importance of understanding the causes of catastrophic forgetting in instruction tuning. Additionally, \cite{ref114} proposes a method to catastrophic forgetting in deep learning, highlighting the importance of addressing this challenge in deep learning.

In conclusion, catastrophic forgetting is a significant challenge in instruction tuning, and addressing it requires a comprehensive approach that combines multiple techniques, including regularization, data augmentation, knowledge distillation, and Bayesian methods. By understanding the causes of catastrophic forgetting and developing effective mitigation strategies, we can improve the performance and reliability of instruction-tuned models, enabling them to learn and adapt to new tasks while retaining their previous capabilities. \cite{ref115} and \cite{ref116} propose methods to preserve ignorance awareness and commonsense knowledge, demonstrating the importance of addressing catastrophic forgetting in LLM fine-tuning and pre-training. Ultimately, addressing catastrophic forgetting is crucial to ensuring the safe and effective deployment of instruction-tuned models in real-world applications.

\subsection{Ethical Concerns and Safety Risks}

The emergence of instruction tuning for large language models (LLMs) has brought about significant advancements in their ability to follow user instructions and generate human-like responses. However, this progress also raises important ethical concerns and safety risks that need to be addressed, particularly in the context of catastrophic forgetting, which can exacerbate these issues. One of the primary concerns is the potential for LLMs to generate harmful or unethical content, which can have severe consequences in real-world applications \cite{ref117}. For instance, if an LLM is instructed to provide information on a sensitive topic, it may inadvertently promote harmful ideologies or provide inaccurate information, which can be detrimental to individuals or society as a whole \cite{ref118}.

The risk of LLMs being used for malicious purposes, such as generating fake news, propaganda, or disinformation, is also a significant concern \cite{ref119}. This can be particularly problematic in situations where LLMs are used to generate content that is intended to influence public opinion or shape cultural narratives. Furthermore, the ability of LLMs to generate convincing and engaging content can make it difficult for users to distinguish between genuine and fake information, which can exacerbate the spread of misinformation \cite{ref120}. In high-stakes applications, such as healthcare or finance, the use of LLMs can also raise concerns about data privacy and security, particularly if sensitive information is shared with the model \cite{ref121}.

To address these ethical concerns and safety risks, it is essential to develop and implement effective mitigation strategies. One approach is to use safety fine-tuning methods, which involve fine-tuning LLMs on datasets that are specifically designed to promote safe and ethical behavior \cite{ref122}. Another approach is to use techniques such as input validation and output filtering to detect and prevent the generation of harmful or unethical content \cite{ref123}. Additionally, it is crucial to develop and implement robust evaluation metrics and testing protocols to ensure that LLMs are safe and effective in real-world applications \cite{ref20}. By prioritizing these concerns and developing effective solutions, we can ensure that instruction-tuned LLMs are developed and deployed in a responsible and safe manner, which is critical for maintaining trust and preventing potential harm.

In conclusion, the development and deployment of instruction-tuned LLMs raise important ethical concerns and safety risks that need to be addressed. These concerns include the potential for LLMs to generate harmful or unethical content, the risk of LLMs being used for malicious purposes, and the safety risks associated with their use in high-stakes applications. To mitigate these risks, it is essential to develop and implement effective mitigation strategies, such as safety fine-tuning methods, input validation and output filtering, and robust evaluation metrics and testing protocols. Furthermore, there is a need for greater transparency and accountability in the development and deployment of LLMs, as well as effective mechanisms for reporting and addressing safety concerns and ethical issues related to LLMs \cite{ref124}. By addressing these concerns and developing effective solutions, we can ensure that instruction-tuned LLMs are developed and deployed in a responsible and safe manner \cite{ref125}.

\section{Comparison of Instruction Tuning Methods}

\subsection{Categorization and Analysis of Instruction Tuning Methods}

The field of instruction tuning has witnessed significant growth in recent years, with various methods being proposed to improve the performance of large language models (LLMs) on diverse tasks. As instruction tuning continues to evolve, it is essential to categorize and analyze different instruction tuning methods, focusing on their performance, efficiency, and applicability to various tasks and domains. This analysis will provide a comprehensive understanding of the strengths and weaknesses of each method, enabling informed decisions about which method to use in a given scenario.

One of the primary categories of instruction tuning methods is supervised fine-tuning, which involves training LLMs on labeled datasets with specific instructions \cite{ref24}. This approach has been shown to be effective in improving the performance of LLMs on tasks such as natural language understanding and generation \cite{ref126}. However, supervised fine-tuning requires large amounts of labeled data, which can be time-consuming and expensive to obtain. In contrast, reinforcement learning from human feedback, which involves training LLMs using rewards or penalties based on human evaluations \cite{ref92}, can be effective in improving the performance of LLMs on tasks such as text generation and dialogue systems \cite{ref127}.

In addition to supervised fine-tuning and reinforcement learning from human feedback, there are other categories of instruction tuning methods, such as prompt engineering and optimization techniques \cite{ref35}. These methods involve designing and optimizing prompts to elicit specific responses from LLMs, and have been shown to be effective in improving the performance of LLMs on tasks such as natural language understanding and generation \cite{ref38}. When analyzing the performance of different instruction tuning methods, it is essential to consider the efficiency and applicability of each method to various tasks and domains. For example, supervised fine-tuning can be computationally expensive and requires large amounts of labeled data, while reinforcement learning from human feedback can be challenging to implement and requires large amounts of human feedback \cite{ref128}.

Instruction tuning methods can be categorized into two main groups: task-specific and task-agnostic methods. Task-specific methods are designed for specific tasks, such as natural language understanding or text generation, and have been shown to be effective in improving the performance of LLMs on these tasks \cite{ref31}. Task-agnostic methods, on the other hand, are designed to be applicable to a wide range of tasks and domains, and have been shown to be effective in improving the performance of LLMs on tasks such as few-shot learning and transfer learning \cite{ref129}. The choice of instruction tuning method depends on the specific task and domain, as well as the available resources and computational budget. By understanding the strengths and weaknesses of each method, researchers and practitioners can make informed decisions about which method to use in a given scenario.

Furthermore, recent studies have explored the use of instruction tuning methods in multimodal settings, such as vision-language models \cite{ref130}. These models have been shown to be effective in improving the performance of LLMs on tasks such as image captioning and visual question answering \cite{ref68}. However, instruction tuning methods in multimodal settings can be challenging to implement, as they require large amounts of labeled data and can be sensitive to the quality of the data. Additionally, instruction tuning methods have also been explored in continual learning scenarios, where the model is trained on a sequence of tasks and must adapt to new tasks and data \cite{ref36}. These methods have been shown to be effective in improving the performance of LLMs on tasks such as few-shot learning and transfer learning \cite{ref131}.

In conclusion, instruction tuning methods are a crucial component of LLMs, and the choice of method depends on the specific task and domain, as well as the available resources and computational budget. As the field of instruction tuning continues to evolve, it is likely to play an increasingly important role in the development of more general, efficient, and safe LLMs. Future research should focus on exploring new instruction tuning methods, such as multimodal and continual learning scenarios, and on developing more efficient and applicable methods for improving the performance of LLMs on diverse tasks and domains \cite{ref32}. The development of more advanced instruction tuning methods will be critical in unlocking the full potential of LLMs and enabling them to be used in a wide range of applications.

\subsection{Recent Advances and Future Directions in Instruction Tuning}

Recent advances in instruction tuning have led to significant improvements in the performance of large language models (LLMs) on various tasks, building upon the foundation established by previous research in the field. As discussed earlier, instruction tuning methods can be categorized into task-specific and task-agnostic methods, each with its strengths and weaknesses. One of the key areas of research has been the development of more efficient and effective methods for instruction tuning, such as the novel data mixture strategy proposed in \cite{ref35}, which uses a submodular function to assign importance scores to tasks and determine the mixture weights. This approach has been shown to significantly outperform traditional methods, demonstrating the potential for improved instruction tuning techniques.

The exploration of different instruction tuning methods has also been a major focus of research, with studies such as \cite{ref12} demonstrating the effectiveness of training with mixed prompt settings, including zero-shot, few-shot, and chain-of-thought prompts. Additionally, \cite{ref24} introduces a simple yet effective task selection method that leverages instruction information alone to identify relevant tasks, optimizing instruction tuning for specific tasks. These advances have paved the way for further research into the applications and limitations of instruction tuning.

The use of instruction tuning has also been explored in multimodal settings, such as vision-language instruction tuning, with \cite{ref11} providing a systematic review of the latest vision-language instruction tuning settings and corresponding datasets in multimodal large language models. The authors identify the characteristics that high-quality vision-language instruction tuning data should possess and conduct extensive experiments to verify their positive impact on the performance of tuned multimodal large language models. Furthermore, researchers have been investigating the impact of instruction tuning on the consistency and reliability of LLMs, with \cite{ref33} comparing the consistency of instruction-tuned LLaMA models to the original LLaMA-7B model and showing that instruction-tuned models become more consistent, both in terms of their representations and their predictions in zero-shot and downstream tasks.

In terms of future directions, one area of research that holds great promise is the development of more advanced instruction tuning methods that can adapt to changing task distributions and user preferences. \cite{ref132} proposes an active instruction tuning method that identifies informative tasks and then actively tunes the model on the selected tasks, achieving better out-of-distribution generalization with fewer training tasks. Another area of research that is likely to gain significant attention in the future is the development of instruction tuning methods that can handle multimodal inputs and outputs, with \cite{ref133} introducing a multimodal instruction tuning benchmark dataset that consists of 62 diverse multimodal tasks in a unified seq-to-seq format.

Additionally, researchers are exploring the use of instruction tuning in real-world applications, such as code generation and program synthesis, with \cite{ref134} proposing a novel fine-tuning technique that enhances LLMs' code generation by extracting multi-granularity differences between correct and incorrect yet similar implementations and dynamically prioritizing those segments during training. Finally, there is a growing interest in developing more explainable and transparent instruction tuning methods that can provide insights into the decision-making processes of LLMs, with \cite{ref2} proposing a framework for interpreting the behavior shift in LLMs after instruction tuning.

As instruction tuning continues to evolve, it is likely to play an increasingly important role in the development of more general, efficient, and safe LLMs. The sample efficiency of instruction-tuned models, as demonstrated in \cite{ref19}, and the inherent instructability of pre-trained language models, highlighted in \cite{ref29}, can be leveraged to develop more effective instruction tuning methods. These advances will likely have a significant impact on the field, enabling the creation of more powerful and flexible LLMs that can be applied to a wide range of tasks and domains, and paving the way for future research into the applications and limitations of instruction tuning.

\section{Future Directions and Open Research Questions}

\subsection{Emerging Trends and Applications}

Emerging trends and applications of instruction tuning are transforming the landscape of natural language processing and its potential impact on various industries and societal applications. One significant trend is the development of more efficient and effective methods for instruction tuning, such as the use of smaller language models \cite{ref135} and the application of mixup-based recipes \cite{ref30}. These advancements have the potential to make instruction tuning more accessible and affordable for a wider range of applications.

The integration of instruction tuning with other techniques, such as neurosymbolic AI \cite{ref136} and multimodal learning \cite{ref17}, is another emerging trend. This integration has the potential to enable more sophisticated and human-like language understanding and generation capabilities. For instance, the use of neurosymbolic AI can enhance the instructability of large language models by providing a more structured and interpretable representation of language \cite{ref136}.

Instruction tuning is also being applied in various industries, such as education \cite{ref137}, healthcare \cite{ref128}, and finance. In education, instruction tuning can be used to develop more effective and personalized learning systems \cite{ref137}. In healthcare, instruction tuning can be used to improve the accuracy and reliability of medical language understanding and generation systems \cite{ref128}. In finance, instruction tuning can be used to develop more sophisticated and human-like language understanding and generation systems for financial analysis and decision-making.

Furthermore, instruction tuning has the potential to impact societal applications, such as improving language accessibility for people with disabilities \cite{ref138} and enhancing language understanding and generation capabilities for low-resource languages \cite{ref17}. For instance, the use of instruction tuning can enable the development of more accurate and reliable language translation systems for low-resource languages \cite{ref17}.

The emergence of large language models has also led to the development of new applications, such as chatbots and virtual assistants \cite{ref139}. These applications have the potential to revolutionize the way we interact with technology and access information. However, they also raise important questions about the potential risks and challenges associated with the use of large language models, such as the potential for bias and misinformation \cite{ref96}.

To address these challenges, researchers are exploring new methods for instruction tuning, such as the use of human curriculum \cite{ref137} and the development of more robust and reliable evaluation metrics \cite{ref127}. These advancements have the potential to enable more effective and efficient instruction tuning, and to mitigate the potential risks and challenges associated with the use of large language models.

In conclusion, emerging trends and applications of instruction tuning are transforming the landscape of natural language processing and its potential impact on various industries and societal applications. The development of more efficient and effective methods for instruction tuning, the integration of instruction tuning with other techniques, and the expansion of instruction tuning to various industries and societal applications are all contributing to the growth and development of this field. However, important questions and challenges remain, and researchers must continue to explore new methods and approaches to address these challenges and to ensure that the benefits of instruction tuning are realized.

The future of instruction tuning holds much promise, with potential applications in areas such as education, healthcare, finance, and more. As researchers continue to explore and develop new methods and techniques, we can expect to see even more sophisticated and human-like language understanding and generation capabilities. The use of instruction tuning has the potential to improve language accessibility for people with disabilities, and to promote greater inclusivity and diversity.

By exploring new methods and approaches, such as the use of human curriculum \cite{ref137} and the development of more robust and reliable evaluation metrics \cite{ref127}, researchers can enable more effective and efficient instruction tuning, and mitigate the potential risks and challenges associated with the use of large language models. The future of instruction tuning is bright, and researchers must continue to explore new methods and approaches to address the challenges and risks associated with the use of large language models. The potential impact of instruction tuning on various industries and societal applications is significant, and researchers must continue to stay up-to-date with the latest developments to ensure that the benefits of instruction tuning are realized.

\subsection{Open Research Questions in Instruction Tuning}

Open research questions in instruction tuning are numerous and varied, reflecting the complexity and multidisciplinary nature of this field. As instruction tuning continues to evolve and transform the landscape of natural language processing, it is essential to identify and address the key challenges and opportunities that lie ahead. One of the significant areas that require further investigation is multimodal instruction tuning, which involves fine-tuning large language models to understand and respond to instructions that include multiple forms of input, such as text, images, and audio \cite{ref133}. This area is crucial for developing models that can interact with humans in a more natural and intuitive way, similar to how humans communicate with each other. However, creating datasets and models that can effectively handle multimodal instructions is a challenging task, especially when considering the need for large, diverse datasets to train these models \cite{ref140}.

The application of instruction tuning in edge AI scenarios is another critical area of research \cite{ref141}. Edge AI refers to the deployment of AI models on edge devices, such as smartphones, smart home devices, and autonomous vehicles, to enable real-time processing and decision-making. Instruction tuning can play a vital role in edge AI by allowing models to be fine-tuned for specific tasks and environments, thereby improving their performance and efficiency. However, edge devices often have limited computational resources and memory, which poses significant challenges for instruction tuning \cite{ref142}. To address this challenge, researchers need to develop methods that can efficiently fine-tune models on edge devices without requiring significant computational resources or large amounts of data.

In addition to these areas, the integration of instruction tuning with other AI techniques, such as reinforcement learning and transfer learning, is another area that requires further exploration \cite{ref143}. Reinforcement learning can be used to fine-tune models based on rewards or penalties received from the environment, while transfer learning enables models to leverage knowledge gained from one task to improve performance on another related task. Combining these techniques with instruction tuning could lead to more robust and adaptable models that can learn from a variety of sources and improve over time. Furthermore, the development of more comprehensive and nuanced evaluation metrics is essential for advancing the field and ensuring that models are developed and deployed in a way that maximizes their potential benefits while minimizing their risks \cite{ref144}.

The ethical and societal implications of instruction tuning and its applications must also be carefully considered \cite{ref145}. As models become more advanced and ubiquitous, there is a growing need for frameworks and guidelines that ensure their development and deployment are transparent, accountable, and aligned with human values. This includes addressing issues related to bias, privacy, and job displacement, as well as ensuring that the benefits of instruction tuning are equitably distributed across different segments of society. The potential of instruction tuning to enhance education and learning is another area that warrants further investigation \cite{ref146}. By developing models that can provide personalized instruction and feedback, instruction tuning could help address some of the longstanding challenges in education, such as unequal access to quality teaching and the need for more effective learning strategies.

Finally, the application of instruction tuning in real-world domains, such as healthcare, finance, and transportation, is an area that requires further research and development \cite{ref147}. In these domains, instruction tuning could be used to develop models that can provide expert-level decision-making and support, leading to improved outcomes and efficiency. However, deploying instruction-tuned models in these domains will require careful consideration of the regulatory, ethical, and societal implications, as well as the development of models that are robust, reliable, and transparent. By addressing these challenges and opportunities, we can unlock the full potential of instruction tuning to drive positive change and improvement in a wide range of domains, and pave the way for the development of more advanced and sophisticated language models that can interact with humans in a more natural and intuitive way.

\subsection{Theoretical Foundations and Future Challenges}

Theoretical foundations of instruction tuning are rooted in the concept of aligning large language models (LLMs) with human preferences and intentions, which is crucial for developing models that can interact with humans in a more natural and intuitive way. As discussed in the previous section, open research questions in instruction tuning are numerous and varied, reflecting the complexity and multidisciplinary nature of this field. The theoretical underpinnings of instruction tuning can be attributed to the emergence of large language models, which have demonstrated impressive capabilities in natural language processing tasks \cite{ref33}. However, the theoretical foundations of instruction tuning also raise several challenges, including the need for high-quality instruction data, the risk of catastrophic forgetting, and the potential for biases in the instruction datasets \cite{ref29}.

One of the key challenges in instruction tuning is the need for high-quality instruction data, which has a significant impact on the performance of the LLM, and low-quality data can lead to suboptimal results \cite{ref13}. To address this challenge, researchers have proposed various methods for improving the quality of instruction data, including data selection and filtering techniques \cite{ref9}. These methods aim to identify the most informative and relevant instruction data, which can be used to fine-tune the LLM and improve its performance. Furthermore, the development of high-quality instruction datasets is closely related to the evaluation metrics and benchmarks used to assess the performance of instruction-tuned models \cite{ref144}, highlighting the need for more comprehensive and nuanced evaluation metrics.

Another challenge in instruction tuning is the risk of catastrophic forgetting, which occurs when an LLM forgets previously learned tasks or knowledge after being fine-tuned on new instruction data \cite{ref148}. To mitigate this risk, researchers have proposed various techniques, including regularization methods and knowledge distillation \cite{ref96}. These techniques aim to preserve the knowledge and capabilities of the LLM while adapting to new instruction data. Additionally, the integration of instruction tuning with other AI techniques, such as reinforcement learning and transfer learning, can help to develop more robust and adaptable models that can learn from a variety of sources and improve over time \cite{ref143}.

Theoretical foundations of instruction tuning also raise questions about the potential biases in the instruction datasets, which can lead to biased LLMs that perpetuate existing social and cultural biases \cite{ref11}. To address this challenge, researchers have proposed various methods for detecting and mitigating biases in instruction datasets, including data augmentation and debiasing techniques \cite{ref99}. Moreover, the development of instruction-tuned models that are transparent, explainable, and aligned with human values is essential for ensuring that the benefits of instruction tuning are equitably distributed across different segments of society \cite{ref145}.

In addition to these challenges, instruction tuning also raises questions about the interpretability and explainability of LLMs, which is essential for understanding and interpreting their behavior \cite{ref2}. To address this challenge, researchers have proposed various techniques, including attention visualization and feature importance methods \cite{ref20}. Future research directions in instruction tuning include the development of more efficient and effective methods for instruction data selection and filtering \cite{ref16}, as well as the exploration of multimodal instruction tuning, which involves fine-tuning LLMs on multimodal instruction datasets that include text, images, and other modalities \cite{ref133}.

Furthermore, future research directions in instruction tuning include the development of more robust and generalizable instruction tuning methods \cite{ref149}, including the exploration of techniques such as meta-learning and few-shot learning, which can enable LLMs to adapt to new instruction datasets and tasks with minimal additional training data \cite{ref150}. Another direction is the investigation of the theoretical foundations of instruction tuning, including the development of formal frameworks and models for understanding the behavior of LLMs during instruction tuning \cite{ref6}. By addressing these challenges and exploring new research directions, we can develop more robust and generalizable instruction tuning methods that can enable LLMs to learn from a wide range of sources and adapt to different contexts and environments \cite{ref127}. Ultimately, the development of instruction-tuned models that are transparent, explainable, and aligned with human values will be crucial for unlocking the full potential of this technology to drive positive change and improvement in a wide range of domains, as will be discussed in the following section.

\section{Conclusion and Recommendations}

\subsection{Summary of Key Findings}

This subsection provides a comprehensive summary of the key findings from the survey on instruction tuning for large language models, laying the foundation for the subsequent discussion on future research directions and recommendations. The survey highlights the importance of instruction tuning in improving the performance of large language models, particularly in natural language understanding and generation tasks \cite{ref151}.

One of the primary findings of the survey is that instruction tuning can significantly improve the performance of large language models on specific tasks, such as text summarization and question answering \cite{ref152}. The survey also emphasizes the importance of dataset quality and instruction data synthesis in instruction tuning, as high-quality instruction data is essential for effective tuning \cite{ref153}.

The survey also explores the challenges and limitations of instruction tuning, including catastrophic forgetting and knowledge retention \cite{ref154}. Catastrophic forgetting occurs when a model forgets previously learned tasks or knowledge during the tuning process, while knowledge retention refers to the ability of a model to retain knowledge learned during pre-training. Understanding these challenges is crucial for developing more effective instruction tuning methods and addressing the limitations of current approaches.

In addition to the technical aspects of instruction tuning, the survey discusses the importance of considering the social and ethical implications of large language models \cite{ref155}. This consideration is essential for ensuring that instruction tuning is developed and applied in a responsible and beneficial manner.

The survey also highlights the potential applications of instruction tuning in various domains, including education, healthcare, and finance \cite{ref156}. For example, instruction tuning can be used to improve the performance of language models in educational settings, such as in automated grading and feedback systems \cite{ref157}.

Overall, the survey provides a comprehensive overview of the current state of instruction tuning for large language models, highlighting the key findings, challenges, and future directions in the field \cite{ref158}. The survey emphasizes the importance of continued research and development in instruction tuning, particularly in addressing the challenges and limitations of current methods and exploring new applications and domains \cite{ref159}. This overview sets the stage for the subsequent discussion on recommendations for future research and practical applications of instruction tuning.

\subsection{Recommendations for Practitioners and Researchers}

As the field of instruction tuning continues to evolve, it is essential for practitioners and researchers to stay informed about the latest developments and best practices. Based on the comprehensive survey of instruction tuning for large language models, several key takeaways and recommendations can be made for future research and practical applications. Firstly, the importance of data quality and instruction data synthesis \cite{ref32} cannot be overstated, as high-quality instruction data can significantly improve the performance of large language models. Therefore, investing time and resources in creating well-structured and diverse datasets is crucial.

Building on this foundation, researchers should explore the potential of adaptive task balancing \cite{ref14} and dynamic knowledge organization \cite{ref36} to improve the efficiency and effectiveness of instruction tuning. These approaches can help mitigate the challenges of catastrophic forgetting and knowledge retention, enabling large language models to learn from a wide range of tasks and adapt to new situations. Furthermore, the development of novel instruction tuning methods, such as Curriculum Instruction Tuning \cite{ref137} and Reward-Oriented Data Selection \cite{ref160}, can provide new avenues for improving the performance of large language models.

In addition to these technical considerations, practitioners and researchers should also prioritize the evaluation of instruction tuning models using a wide range of metrics and benchmarks \cite{ref161}. This can help ensure that the models are generalizable and effective in real-world applications, and can provide a more comprehensive understanding of their strengths and limitations. By adopting a multifaceted approach to instruction tuning, researchers can unlock the full potential of large language models and drive innovation in areas such as natural language processing, computer vision, and human-computer interaction.

Looking ahead, several areas can be identified as promising and worthy of exploration. The development of more efficient and effective instruction tuning methods, such as those using submodular functions \cite{ref35} and token-wise attention-derived saliency \cite{ref37}, can help reduce the computational costs and improve the performance of large language models. Additionally, the exploration of instruction tuning in multimodal and edge AI applications \cite{ref133} can enable the creation of more versatile and effective large language models, with potential applications in areas such as computer vision, natural language processing, and human-computer interaction.

The investigation of the impact of instruction tuning on the susceptibility of large language models to misinformation \cite{ref96} is also a critical area of research, as it can help identify potential risks and challenges associated with the use of instruction tuning. By exploring these areas and developing more robust and reliable models, researchers can ensure that the benefits of instruction tuning are realized while minimizing its potential drawbacks. Ultimately, the future of instruction tuning holds tremendous promise, and by working together to advance the field, practitioners and researchers can create more effective, efficient, and versatile large language models that can drive innovation and improve outcomes in a wide range of domains.

\subsection{Implications and Potential Impact}

The implications of instruction tuning for the development of more general, efficient, and safe large language models are profound, and they have significant connections to the future research directions and recommendations outlined earlier. As demonstrated by various studies, instruction tuning has the potential to significantly improve the performance of large language models on a wide range of tasks, from natural language understanding and generation to code generation and multimodal tasks \cite{ref11}. This improvement in performance can have a substantial impact on various industries, such as customer service, language translation, and content creation, where large language models are increasingly being used to automate tasks and improve efficiency. Furthermore, the potential of instruction tuning to enhance the performance of large language models is closely related to the development of novel instruction tuning methods, such as Curriculum Instruction Tuning \cite{ref137} and Reward-Oriented Data Selection \cite{ref160}, which can help identify the most informative and relevant instruction data.

Moreover, instruction tuning can also lead to the development of more efficient large language models, which is in line with the recommendation to explore the potential of adaptive task balancing \cite{ref14} and dynamic knowledge organization \cite{ref36}. By fine-tuning models on specific tasks and datasets, researchers can reduce the computational resources required for training and improve the models' ability to generalize to new tasks \cite{ref162}. This can be particularly beneficial for smaller organizations or individuals who may not have access to large amounts of computational resources. Additionally, the use of instruction tuning can also lead to more efficient models in terms of inference time, as demonstrated by \cite{ref9}, which can be critical for real-time applications.

The potential impact of instruction tuning on societal applications is also significant, and it is closely related to the recommendation to investigate the potential of instruction tuning in real-world applications, such as recommender systems \cite{ref163} and natural language processing \cite{ref88}. For instance, instruction-tuned models can be used to improve language understanding and generation for low-resource languages, which can have a positive impact on communities that speak these languages \cite{ref139}. Furthermore, instruction tuning can also be used to develop more safe and reliable large language models, as demonstrated by \cite{ref117}, which can be critical for applications such as healthcare and education.

However, the development of instruction-tuned models also raises important questions about the potential risks and challenges associated with these models, which is in line with the recommendation to investigate the impact of instruction tuning on the susceptibility of large language models to misinformation \cite{ref96}. For example, the use of instruction tuning can lead to models that are more prone to bias and discrimination, as demonstrated by \cite{ref164}, which can have negative consequences for certain groups of people. Additionally, the development of more advanced instruction-tuned models can also lead to job displacement and exacerbate existing social inequalities, as highlighted by \cite{ref2}.

To mitigate these risks, it is essential to develop more transparent and explainable instruction-tuned models, as proposed by \cite{ref165}, and to develop more diverse and representative instruction datasets, as highlighted by \cite{ref166}. This can involve developing new evaluation metrics and methodologies that can assess the performance and safety of instruction-tuned models, as well as developing more robust and reliable models that can generalize to new tasks and datasets. By addressing these challenges and developing more responsible instruction-tuned models, we can unlock the full potential of instruction tuning and create more effective, efficient, and safe large language models.

In conclusion, the implications of instruction tuning for the development of more general, efficient, and safe large language models are significant, and they have important connections to the future research directions and recommendations outlined earlier. While there are potential risks and challenges associated with instruction-tuned models, these can be mitigated by developing more transparent and explainable models, as well as more diverse and representative instruction datasets. As the field of natural language processing continues to evolve, it is likely that instruction tuning will play an increasingly important role in the development of more advanced and reliable large language models, with significant potential impacts on various industries and societal applications \cite{ref32}.


\begin{thebibliography}{166}
\addcontentsline{toc}{section}{References}
\bibitem{ref1} Dynamics of Instruction Fine-Tuning for Chinese Large Language Models. arXiv:2310.19651, 2023. \url{https://arxiv.org/abs/2310.19651}
\bibitem{ref2} From Language Modeling to Instruction Following: Understanding the Behavior Shift in LLMs after Instruction Tuning. arXiv:2310.00492, 2023. \url{https://arxiv.org/abs/2310.00492}
\bibitem{ref3} Evaluating the Zero-shot Robustness of Instruction-tuned Language Models. arXiv:2306.11270, 2023. \url{https://arxiv.org/abs/2306.11270}
\bibitem{ref4} TIM: Teaching Large Language Models to Translate with Comparison. arXiv:2307.04408, 2023. \url{https://arxiv.org/abs/2307.04408}
\bibitem{ref5} Instruction-tuning Aligns LLMs to the Human Brain. arXiv:2312.00575, 2023. \url{https://arxiv.org/abs/2312.00575}
\bibitem{ref6} The Inherent Limits of Pretrained LLMs: The Unexpected Convergence of Instruction Tuning and In-Context Learning Capabilities. arXiv:2501.08716, 2025. \url{https://arxiv.org/abs/2501.08716}
\bibitem{ref7} From Base to Conversational: Japanese Instruction Dataset and Tuning Large Language Models. arXiv:2309.03412, 2023. \url{https://arxiv.org/abs/2309.03412}
\bibitem{ref8} A Survey on Data Selection for LLM Instruction Tuning. arXiv:2402.05123, 2024. \url{https://arxiv.org/abs/2402.05123}
\bibitem{ref9} Superfiltering: Weak-to-Strong Data Filtering for Fast Instruction-Tuning. arXiv:2402.00530, 2024. \url{https://arxiv.org/abs/2402.00530}
\bibitem{ref10} Zero-shot cross-lingual transfer in instruction tuning of large language models. arXiv:2402.14778, 2024. \url{https://arxiv.org/abs/2402.14778}
\bibitem{ref11} Vision-Language Instruction Tuning: A Review and Analysis. arXiv:2311.08172, 2023. \url{https://arxiv.org/abs/2311.08172}
\bibitem{ref12} The Flan Collection: Designing Data and Methods for Effective Instruction Tuning. arXiv:2301.13688, 2023. \url{https://arxiv.org/abs/2301.13688}
\bibitem{ref13} Rethinking the Instruction Quality: LIFT is What You Need. arXiv:2312.11508, 2023. \url{https://arxiv.org/abs/2312.11508}
\bibitem{ref14} Adaptive Task Balancing for Visual Instruction Tuning via Inter-Task Contribution and Intra-Task Difficulty. arXiv:2403.04343, 2024. \url{https://arxiv.org/abs/2403.04343}
\bibitem{ref15} Optimizing Instruction Synthesis: Effective Exploration of Evolutionary Space with Tree Search. arXiv:2410.10392, 2024. \url{https://arxiv.org/abs/2410.10392}
\bibitem{ref16} LESS: Selecting Influential Data for Targeted Instruction Tuning. arXiv:2402.04333, 2024. \url{https://arxiv.org/abs/2402.04333}
\bibitem{ref17} M\$\textasciicircum{}3\$IT: A Large-Scale Dataset towards Multi-Modal Multilingual Instruction Tuning. arXiv:2306.04387, 2023. \url{https://arxiv.org/abs/2306.04387}
\bibitem{ref18} SelfCodeAlign: Self-Alignment for Code Generation. arXiv:2410.24198, 2024. \url{https://arxiv.org/abs/2410.24198}
\bibitem{ref19} Instruction Tuned Models are Quick Learners. arXiv:2306.05539, 2023. \url{https://arxiv.org/abs/2306.05539}
\bibitem{ref20} Revisiting Instruction Fine-tuned Model Evaluation to Guide Industrial Applications. arXiv:2310.14103, 2023. \url{https://arxiv.org/abs/2310.14103}
\bibitem{ref21} What Makes Good Data for Alignment? A Comprehensive Study of Automatic Data Selection in Instruction Tuning. arXiv:2312.15685, 2023. \url{https://arxiv.org/abs/2312.15685}
\bibitem{ref22} Toward Secure Tuning: Mitigating Security Risks from Instruction Fine-Tuning. arXiv:2410.04524, 2024. \url{https://arxiv.org/abs/2410.04524}
\bibitem{ref23} Continual LLaVA: Continual Instruction Tuning in Large Vision-Language Models. arXiv:2411.02564, 2024. \url{https://arxiv.org/abs/2411.02564}
\bibitem{ref24} Instruction Matters: A Simple yet Effective Task Selection for Optimized Instruction Tuning of Specific Tasks. arXiv:2404.16418, 2024. \url{https://arxiv.org/abs/2404.16418}
\bibitem{ref25} Mixture-of-Experts Meets Instruction Tuning:A Winning Combination for Large Language Models. arXiv:2305.14705, 2023. \url{https://arxiv.org/abs/2305.14705}
\bibitem{ref26} Unveiling the Secret Recipe: A Guide For Supervised Fine-Tuning Small LLMs. arXiv:2412.13337, 2024. \url{https://arxiv.org/abs/2412.13337}
\bibitem{ref27} Layer by Layer: Uncovering Where Multi-Task Learning Happens in Instruction-Tuned Large Language Models. arXiv:2410.20008, 2024. \url{https://arxiv.org/abs/2410.20008}
\bibitem{ref28} Panda LLM: Training Data and Evaluation for Open-Sourced Chinese Instruction-Following Large Language Models. arXiv:2305.03025, 2023. \url{https://arxiv.org/abs/2305.03025}
\bibitem{ref29} Revealing the Inherent Instructability of Pre-Trained Language Models. arXiv:2410.02465, 2024. \url{https://arxiv.org/abs/2410.02465}
\bibitem{ref30} SFTMix: Elevating Language Model Instruction Tuning with Mixup Recipe. arXiv:2410.05248, 2024. \url{https://arxiv.org/abs/2410.05248}
\bibitem{ref31} Instruction Tuning With Loss Over Instructions. arXiv:2405.14394, 2024. \url{https://arxiv.org/abs/2405.14394}
\bibitem{ref32} Unleashing the Power of Data Tsunami: A Comprehensive Survey on Data Assessment and Selection for Instruction Tuning of Language Models. arXiv:2408.02085, 2024. \url{https://arxiv.org/abs/2408.02085}
\bibitem{ref33} Does Instruction Tuning Make LLMs More Consistent?. arXiv:2404.15206, 2024. \url{https://arxiv.org/abs/2404.15206}
\bibitem{ref34} MoDS: Model-oriented Data Selection for Instruction Tuning. arXiv:2311.15653, 2023. \url{https://arxiv.org/abs/2311.15653}
\bibitem{ref35} SMART: Submodular Data Mixture Strategy for Instruction Tuning. arXiv:2403.08370, 2024. \url{https://arxiv.org/abs/2403.08370}
\bibitem{ref36} Federated Continual Instruction Tuning. arXiv:2503.12897, 2025. \url{https://arxiv.org/abs/2503.12897}
\bibitem{ref37} TRIM: Token-wise Attention-Derived Saliency for Data-Efficient Instruction Tuning. arXiv:2510.07118, 2025. \url{https://arxiv.org/abs/2510.07118}
\bibitem{ref38} PACIT: Unlocking the Power of Examples for Better In-Context Instruction Tuning. arXiv:2310.00901, 2023. \url{https://arxiv.org/abs/2310.00901}
\bibitem{ref39} IterSelectTune: An Iterative Training Framework for Efficient Instruction-Tuning Data Selection. arXiv:2410.13464, 2024. \url{https://arxiv.org/abs/2410.13464}
\bibitem{ref40} InsCL: A Data-efficient Continual Learning Paradigm for Fine-tuning Large Language Models with Instructions. arXiv:2403.11435, 2024. \url{https://arxiv.org/abs/2403.11435}
\bibitem{ref41} Prototypical Reward Network for Data-Efficient RLHF. arXiv:2406.06606, 2024. \url{https://arxiv.org/abs/2406.06606}
\bibitem{ref42} Interpreting Learned Feedback Patterns in Large Language Models. arXiv:2310.08164, 2023. \url{https://arxiv.org/abs/2310.08164}
\bibitem{ref43} Training a Helpful and Harmless Assistant with Reinforcement Learning from Human Feedback. arXiv:2204.05862, 2022. \url{https://arxiv.org/abs/2204.05862}
\bibitem{ref44} Understanding and Alleviating Memory Consumption in RLHF for LLMs. arXiv:2410.15651, 2024. \url{https://arxiv.org/abs/2410.15651}
\bibitem{ref45} Shattering the Agent-Environment Interface for Fine-Tuning Inclusive Language Models. arXiv:2305.11455, 2023. \url{https://arxiv.org/abs/2305.11455}
\bibitem{ref46} SuperHF: Supervised Iterative Learning from Human Feedback. arXiv:2310.16763, 2023. \url{https://arxiv.org/abs/2310.16763}
\bibitem{ref47} Fine-Grained Human Feedback Gives Better Rewards for Language Model Training. arXiv:2306.01693, 2023. \url{https://arxiv.org/abs/2306.01693}
\bibitem{ref48} Sequence to Sequence Reward Modeling: Improving RLHF by Language Feedback. arXiv:2409.00162, 2024. \url{https://arxiv.org/abs/2409.00162}
\bibitem{ref49} Equilibrate RLHF: Towards Balancing Helpfulness-Safety Trade-off in Large Language Models. arXiv:2502.11555, 2025. \url{https://arxiv.org/abs/2502.11555}
\bibitem{ref50} Evaluating Parameter-Efficient Transfer Learning Approaches on SURE Benchmark for Speech Understanding. arXiv:2303.03267, 2023. \url{https://arxiv.org/abs/2303.03267}
\bibitem{ref51} Fine-tuning Reinforcement Learning Models is Secretly a Forgetting Mitigation Problem. arXiv:2402.02868, 2024. \url{https://arxiv.org/abs/2402.02868}
\bibitem{ref52} Rethinking the Hyperparameters for Fine-tuning. arXiv:2002.11770, 2020. \url{https://arxiv.org/abs/2002.11770}
\bibitem{ref53} Know Where You're Going: Meta-Learning for Parameter-Efficient Fine-Tuning. arXiv:2205.12453, 2022. \url{https://arxiv.org/abs/2205.12453}
\bibitem{ref54} Parameter-Efficient Fine-Tuning of Large Language Models for Unit Test Generation: An Empirical Study. arXiv:2411.02462, 2024. \url{https://arxiv.org/abs/2411.02462}
\bibitem{ref55} Prefix-Tuning+: Modernizing Prefix-Tuning by Decoupling the Prefix from Attention. arXiv:2506.13674, 2025. \url{https://arxiv.org/abs/2506.13674}
\bibitem{ref56} Transfer Learning for Finetuning Large Language Models. arXiv:2411.01195, 2024. \url{https://arxiv.org/abs/2411.01195}
\bibitem{ref57} MEFT: Memory-Efficient Fine-Tuning through Sparse Adapter. arXiv:2406.04984, 2024. \url{https://arxiv.org/abs/2406.04984}
\bibitem{ref58} Gradual Tuning: a better way of Fine Tuning the parameters of a Deep Neural Network. arXiv:1711.10177, 2017. \url{https://arxiv.org/abs/1711.10177}
\bibitem{ref59} TAIA: Large Language Models are Out-of-Distribution Data Learners. arXiv:2405.20192, 2024. \url{https://arxiv.org/abs/2405.20192}
\bibitem{ref60} Parameter-Efficient Fine-Tuning without Introducing New Latency. arXiv:2305.16742, 2023. \url{https://arxiv.org/abs/2305.16742}
\bibitem{ref61} Astraios: Parameter-Efficient Instruction Tuning Code Large Language Models. arXiv:2401.00788, 2024. \url{https://arxiv.org/abs/2401.00788}
\bibitem{ref62} P-Tuning v2: Prompt Tuning Can Be Comparable to Fine-tuning Universally Across Scales and Tasks. arXiv:2110.07602, 2021. \url{https://arxiv.org/abs/2110.07602}
\bibitem{ref63} QPO: Query-dependent Prompt Optimization via Multi-Loop Offline Reinforcement Learning. arXiv:2408.10504, 2024. \url{https://arxiv.org/abs/2408.10504}
\bibitem{ref64} TEMPERA: Test-Time Prompting via Reinforcement Learning. arXiv:2211.11890, 2022. \url{https://arxiv.org/abs/2211.11890}
\bibitem{ref65} GrIPS: Gradient-free, Edit-based Instruction Search for Prompting Large Language Models. arXiv:2203.07281, 2022. \url{https://arxiv.org/abs/2203.07281}
\bibitem{ref66} M\$\textasciicircum{}2\$PT: Multimodal Prompt Tuning for Zero-shot Instruction Learning. arXiv:2409.15657, 2024. \url{https://arxiv.org/abs/2409.15657}
\bibitem{ref67} Task Facet Learning: A Structured Approach to Prompt Optimization. arXiv:2406.10504, 2024. \url{https://arxiv.org/abs/2406.10504}
\bibitem{ref68} Re-Imagining Multimodal Instruction Tuning: A Representation View. arXiv:2503.00723, 2025. \url{https://arxiv.org/abs/2503.00723}
\bibitem{ref69} AMPO: Automatic Multi-Branched Prompt Optimization. arXiv:2410.08696, 2024. \url{https://arxiv.org/abs/2410.08696}
\bibitem{ref70} Do Models Really Learn to Follow Instructions? An Empirical Study of Instruction Tuning. arXiv:2305.11383, 2023. \url{https://arxiv.org/abs/2305.11383}
\bibitem{ref71} Finetuned Language Models Are Zero-Shot Learners. arXiv:2109.01652, 2021. \url{https://arxiv.org/abs/2109.01652}
\bibitem{ref72} Raw Text is All you Need: Knowledge-intensive Multi-turn Instruction Tuning for Large Language Model. arXiv:2407.03040, 2024. \url{https://arxiv.org/abs/2407.03040}
\bibitem{ref73} X-Eval: Generalizable Multi-aspect Text Evaluation via Augmented Instruction Tuning with Auxiliary Evaluation Aspects. arXiv:2311.08788, 2023. \url{https://arxiv.org/abs/2311.08788}
\bibitem{ref74} MIMIC-IT: Multi-Modal In-Context Instruction Tuning. arXiv:2306.05425, 2023. \url{https://arxiv.org/abs/2306.05425}
\bibitem{ref75} BioInstruct: Instruction Tuning of Large Language Models for Biomedical Natural Language Processing. arXiv:2310.19975, 2023. \url{https://arxiv.org/abs/2310.19975}
\bibitem{ref76} INTERS: Unlocking the Power of Large Language Models in Search with Instruction Tuning. arXiv:2401.06532, 2024. \url{https://arxiv.org/abs/2401.06532}
\bibitem{ref77} Chain-of-Instructions: Compositional Instruction Tuning on Large Language Models. arXiv:2402.11532, 2024. \url{https://arxiv.org/abs/2402.11532}
\bibitem{ref78} SQATIN: Supervised Instruction Tuning Meets Question Answering for Improved Dialogue NLU. arXiv:2311.09502, 2023. \url{https://arxiv.org/abs/2311.09502}
\bibitem{ref79} SemScore: Automated Evaluation of Instruction-Tuned LLMs based on Semantic Textual Similarity. arXiv:2401.17072, 2024. \url{https://arxiv.org/abs/2401.17072}
\bibitem{ref80} How Far Can Camels Go? Exploring the State of Instruction Tuning on Open Resources. arXiv:2306.04751, 2023. \url{https://arxiv.org/abs/2306.04751}
\bibitem{ref81} PandaLM: An Automatic Evaluation Benchmark for LLM Instruction Tuning Optimization. arXiv:2306.05087, 2023. \url{https://arxiv.org/abs/2306.05087}
\bibitem{ref82} Evaluating the Robustness to Instructions of Large Language Models. arXiv:2308.14306, 2023. \url{https://arxiv.org/abs/2308.14306}
\bibitem{ref83} Toward the Evaluation of Large Language Models Considering Score Variance across Instruction Templates. arXiv:2408.12263, 2024. \url{https://arxiv.org/abs/2408.12263}
\bibitem{ref84} FollowEval: A Multi-Dimensional Benchmark for Assessing the Instruction-Following Capability of Large Language Models. arXiv:2311.09829, 2023. \url{https://arxiv.org/abs/2311.09829}
\bibitem{ref85} Evaluating Correctness and Faithfulness of Instruction-Following Models for Question Answering. arXiv:2307.16877, 2023. \url{https://arxiv.org/abs/2307.16877}
\bibitem{ref86} DecompEval: Evaluating Generated Texts as Unsupervised Decomposed Question Answering. arXiv:2307.06869, 2023. \url{https://arxiv.org/abs/2307.06869}
\bibitem{ref87} TIME: Temporal-Sensitive Multi-Dimensional Instruction Tuning and Robust Benchmarking for Video-LLMs. arXiv:2503.09994, 2025. \url{https://arxiv.org/abs/2503.09994}
\bibitem{ref88} From Symbolic Tasks to Code Generation: Diversification Yields Better Task Performers. arXiv:2405.19787, 2024. \url{https://arxiv.org/abs/2405.19787}
\bibitem{ref89} Scalable Fine-tuning from Multiple Data Sources: A First-Order Approximation Approach. arXiv:2409.19458, 2024. \url{https://arxiv.org/abs/2409.19458}
\bibitem{ref90} OPT-IML: Scaling Language Model Instruction Meta Learning through the Lens of Generalization. arXiv:2212.12017, 2022. \url{https://arxiv.org/abs/2212.12017}
\bibitem{ref91} Maybe Only 0.5\% Data is Needed: A Preliminary Exploration of Low Training Data Instruction Tuning. arXiv:2305.09246, 2023. \url{https://arxiv.org/abs/2305.09246}
\bibitem{ref92} Harnessing the Power of David against Goliath: Exploring Instruction Data Generation without Using Closed-Source Models. arXiv:2308.12711, 2023. \url{https://arxiv.org/abs/2308.12711}
\bibitem{ref93} Large-Scale Data Selection for Instruction Tuning. arXiv:2503.01807, 2025. \url{https://arxiv.org/abs/2503.01807}
\bibitem{ref94} MMInstruct: a high-quality multi-modal instruction tuning dataset with extensive diversity. arXiv:2407.15838, 2024. \url{https://arxiv.org/abs/2407.15838}
\bibitem{ref95} Less is More: High-value Data Selection for Visual Instruction Tuning. arXiv:2403.09559, 2024. \url{https://arxiv.org/abs/2403.09559}
\bibitem{ref96} Exploring the Impact of Instruction-Tuning on LLM's Susceptibility to Misinformation. arXiv:2507.18203, 2025. \url{https://arxiv.org/abs/2507.18203}
\bibitem{ref97} Data Quality Control in Federated Instruction-tuning of Large Language Models. arXiv:2410.11540, 2024. \url{https://arxiv.org/abs/2410.11540}
\bibitem{ref98} EasyInstruct: An Easy-to-use Instruction Processing Framework for Large Language Models. arXiv:2402.03049, 2024. \url{https://arxiv.org/abs/2402.03049}
\bibitem{ref99} Federated Data-Efficient Instruction Tuning for Large Language Models. arXiv:2410.10926, 2024. \url{https://arxiv.org/abs/2410.10926}
\bibitem{ref100} Personalized Federated Instruction Tuning via Neural Architecture Search. arXiv:2402.16919, 2024. \url{https://arxiv.org/abs/2402.16919}
\bibitem{ref101} An Empirical Study of Catastrophic Forgetting in Large Language Models During Continual Fine-tuning. arXiv:2308.08747, 2023. \url{https://arxiv.org/abs/2308.08747}
\bibitem{ref102} Refine Large Language Model Fine-tuning via Instruction Vector. arXiv:2406.12227, 2024. \url{https://arxiv.org/abs/2406.12227}
\bibitem{ref103} SwitchCIT: Switching for Continual Instruction Tuning. arXiv:2407.11780, 2024. \url{https://arxiv.org/abs/2407.11780}
\bibitem{ref104} Chained Tuning Leads to Biased Forgetting. arXiv:2412.16469, 2024. \url{https://arxiv.org/abs/2412.16469}
\bibitem{ref105} Context-Free Synthetic Data Mitigates Forgetting. arXiv:2505.13811, 2025. \url{https://arxiv.org/abs/2505.13811}
\bibitem{ref106} HFT: Half Fine-Tuning for Large Language Models. arXiv:2404.18466, 2024. \url{https://arxiv.org/abs/2404.18466}
\bibitem{ref107} MCF-VC: Mitigate Catastrophic Forgetting in Class-Incremental Learning for Multimodal Video Captioning. arXiv:2402.17680, 2024. \url{https://arxiv.org/abs/2402.17680}
\bibitem{ref108} ClaRe: Practical Class Incremental Learning By Remembering Previous Class Representations. arXiv:2103.15486, 2021. \url{https://arxiv.org/abs/2103.15486}
\bibitem{ref109} CODA-Prompt: COntinual Decomposed Attention-based Prompting for Rehearsal-Free Continual Learning. arXiv:2211.13218, 2022. \url{https://arxiv.org/abs/2211.13218}
\bibitem{ref110} SEFE: Superficial and Essential Forgetting Eliminator for Multimodal Continual Instruction Tuning. arXiv:2505.02486, 2025. \url{https://arxiv.org/abs/2505.02486}
\bibitem{ref111} Bayesian Parameter-Efficient Fine-Tuning for Overcoming Catastrophic Forgetting. arXiv:2402.12220, 2024. \url{https://arxiv.org/abs/2402.12220}
\bibitem{ref112} Dynamic Orthogonal Continual Fine-tuning for Mitigating Catastrophic Forgettings. arXiv:2509.23893, 2025. \url{https://arxiv.org/abs/2509.23893}
\bibitem{ref113} Anatomy of Catastrophic Forgetting: Hidden Representations and Task Semantics. arXiv:2007.07400, 2020. \url{https://arxiv.org/abs/2007.07400}
\bibitem{ref114} Catastrophic Forgetting in Deep Learning: A Comprehensive Taxonomy. arXiv:2312.10549, 2023. \url{https://arxiv.org/abs/2312.10549}
\bibitem{ref115} Don't Make It Up: Preserving Ignorance Awareness in LLM Fine-Tuning. arXiv:2506.14387, 2025. \url{https://arxiv.org/abs/2506.14387}
\bibitem{ref116} Preserving Commonsense Knowledge from Pre-trained Language Models via Causal Inference. arXiv:2306.10790, 2023. \url{https://arxiv.org/abs/2306.10790}
\bibitem{ref117} Safety-Tuned LLaMAs: Lessons From Improving the Safety of Large Language Models that Follow Instructions. arXiv:2309.07875, 2023. \url{https://arxiv.org/abs/2309.07875}
\bibitem{ref118} How (un)ethical are instruction-centric responses of LLMs? Unveiling the vulnerabilities of safety guardrails to harmful queries. arXiv:2402.15302, 2024. \url{https://arxiv.org/abs/2402.15302}
\bibitem{ref119} Backdooring Instruction-Tuned Large Language Models with Virtual Prompt Injection. arXiv:2307.16888, 2023. \url{https://arxiv.org/abs/2307.16888}
\bibitem{ref120} Mitigating Covertly Unsafe Text within Natural Language Systems. arXiv:2210.09306, 2022. \url{https://arxiv.org/abs/2210.09306}
\bibitem{ref121} Safe Generative Chats in a WhatsApp Intelligent Tutoring System. arXiv:2407.04915, 2024. \url{https://arxiv.org/abs/2407.04915}
\bibitem{ref122} Safety Fine-Tuning at (Almost) No Cost: A Baseline for Vision Large Language Models. arXiv:2402.02207, 2024. \url{https://arxiv.org/abs/2402.02207}
\bibitem{ref123} Safety Arithmetic: A Framework for Test-time Safety Alignment of Language Models by Steering Parameters and Activations. arXiv:2406.11801, 2024. \url{https://arxiv.org/abs/2406.11801}
\bibitem{ref124} LLM Can be a Dangerous Persuader: Empirical Study of Persuasion Safety in Large Language Models. arXiv:2504.10430, 2025. \url{https://arxiv.org/abs/2504.10430}
\bibitem{ref125} Course-Correction: Safety Alignment Using Synthetic Preferences. arXiv:2407.16637, 2024. \url{https://arxiv.org/abs/2407.16637}
\bibitem{ref126} Demystifying Instruction Mixing for Fine-tuning Large Language Models. arXiv:2312.10793, 2023. \url{https://arxiv.org/abs/2312.10793}
\bibitem{ref127} InstructDial: Improving Zero and Few-shot Generalization in Dialogue through Instruction Tuning. arXiv:2205.12673, 2022. \url{https://arxiv.org/abs/2205.12673}
\bibitem{ref128} Fine-Tuning on Noisy Instructions: Effects on Generalization and Performance. arXiv:2510.03528, 2025. \url{https://arxiv.org/abs/2510.03528}
\bibitem{ref129} DELIFT: Data Efficient Language model Instruction Fine Tuning. arXiv:2411.04425, 2024. \url{https://arxiv.org/abs/2411.04425}
\bibitem{ref130} Visual Instruction Tuning towards General-Purpose Multimodal Model: A Survey. arXiv:2312.16602, 2023. \url{https://arxiv.org/abs/2312.16602}
\bibitem{ref131} CoIN: A Benchmark of Continual Instruction tuNing for Multimodel Large Language Model. arXiv:2403.08350, 2024. \url{https://arxiv.org/abs/2403.08350}
\bibitem{ref132} Active Instruction Tuning: Improving Cross-Task Generalization by Training on Prompt Sensitive Tasks. arXiv:2311.00288, 2023. \url{https://arxiv.org/abs/2311.00288}
\bibitem{ref133} MultiInstruct: Improving Multi-Modal Zero-Shot Learning via Instruction Tuning. arXiv:2212.10773, 2022. \url{https://arxiv.org/abs/2212.10773}
\bibitem{ref134} FGIT: Fault-Guided Fine-Tuning for Code Generation. arXiv:2503.16913, 2025. \url{https://arxiv.org/abs/2503.16913}
\bibitem{ref135} Smaller Language Models Are Better Instruction Evolvers. arXiv:2412.11231, 2024. \url{https://arxiv.org/abs/2412.11231}
\bibitem{ref136} Neurosymbolic AI for Enhancing Instructability in Generative AI. arXiv:2407.18722, 2024. \url{https://arxiv.org/abs/2407.18722}
\bibitem{ref137} Instruction Tuning with Human Curriculum. arXiv:2310.09518, 2023. \url{https://arxiv.org/abs/2310.09518}
\bibitem{ref138} Instruction Tuning for Secure Code Generation. arXiv:2402.09497, 2024. \url{https://arxiv.org/abs/2402.09497}
\bibitem{ref139} Chinese Open Instruction Generalist: A Preliminary Release. arXiv:2304.07987, 2023. \url{https://arxiv.org/abs/2304.07987}
\bibitem{ref140} MMDU: A Multi-Turn Multi-Image Dialog Understanding Benchmark and Instruction-Tuning Dataset for LVLMs. arXiv:2406.11833, 2024. \url{https://arxiv.org/abs/2406.11833}
\bibitem{ref141} Green Edge AI: A Contemporary Survey. arXiv:2312.00333, 2023. \url{https://arxiv.org/abs/2312.00333}
\bibitem{ref142} NVCiM-PT: An NVCiM-assisted Prompt Tuning Framework for Edge LLMs. arXiv:2411.08244, 2024. \url{https://arxiv.org/abs/2411.08244}
\bibitem{ref143} When Parameter-efficient Tuning Meets General-purpose Vision-language Models. arXiv:2312.12458, 2023. \url{https://arxiv.org/abs/2312.12458}
\bibitem{ref144} ProBench: Judging Multimodal Foundation Models on Open-ended Multi-domain Expert Tasks. arXiv:2503.06885, 2025. \url{https://arxiv.org/abs/2503.06885}
\bibitem{ref145} AI Governance in the Context of the EU AI Act: A Bibliometric and Literature Review Approach. arXiv:2502.03468, 2025. \url{https://arxiv.org/abs/2502.03468}
\bibitem{ref146} New Era of Artificial Intelligence in Education: Towards a Sustainable Multifaceted Revolution. arXiv:2305.18303, 2023. \url{https://arxiv.org/abs/2305.18303}
\bibitem{ref147} Advancements in Mobile Edge Computing and Open RAN: Leveraging Artificial Intelligence and Machine Learning for Wireless Systems. arXiv:2502.02886, 2025. \url{https://arxiv.org/abs/2502.02886}
\bibitem{ref148} Anchoring Refusal Direction: Mitigating Safety Risks in Tuning via Projection Constraint. arXiv:2509.06795, 2025. \url{https://arxiv.org/abs/2509.06795}
\bibitem{ref149} Dual Instruction Tuning with Large Language Models for Mathematical Reasoning. arXiv:2403.18295, 2024. \url{https://arxiv.org/abs/2403.18295}
\bibitem{ref150} One-Shot Learning as Instruction Data Prospector for Large Language Models. arXiv:2312.10302, 2023. \url{https://arxiv.org/abs/2312.10302}
\bibitem{ref151} An Empirical Survey on Long Document Summarization: Datasets, Models, and Metrics. arXiv:2207.00939, 2022. \url{https://arxiv.org/abs/2207.00939}
\bibitem{ref152} What Makes a Good and Useful Summary? Incorporating Users in Automatic Summarization Research. arXiv:2012.07619, 2020. \url{https://arxiv.org/abs/2012.07619}
\bibitem{ref153} Ensuring anonymity in survey panel research. arXiv:1404.1808, 2014. \url{https://arxiv.org/abs/1404.1808}
\bibitem{ref154} Subfield Effects on the Core of Coauthors. arXiv:1306.0453, 2013. \url{https://arxiv.org/abs/1306.0453}
\bibitem{ref155} Economic Diversification and Social Progress in the GCC Countries: A Study on the Transition from Oil-Dependency to Knowledge-Based Economies. arXiv:2410.21505, 2024. \url{https://arxiv.org/abs/2410.21505}
\bibitem{ref156} Influencing factors of Twitter mentions of scientific papers. arXiv:2210.11852, 2022. \url{https://arxiv.org/abs/2210.11852}
\bibitem{ref157} Machine Learning Predicts Upper Secondary Education Dropout as Early as the End of Primary School. arXiv:2403.14663, 2024. \url{https://arxiv.org/abs/2403.14663}
\bibitem{ref158} A Systematic Survey of Text Summarization: From Statistical Methods to Large Language Models. arXiv:2406.11289, 2024. \url{https://arxiv.org/abs/2406.11289}
\bibitem{ref159} Determining Research Priorities for Astronomy Using Machine Learning. arXiv:2203.00713, 2022. \url{https://arxiv.org/abs/2203.00713}
\bibitem{ref160} ROSE: A Reward-Oriented Data Selection Framework for LLM Task-Specific Instruction Tuning. arXiv:2412.00631, 2024. \url{https://arxiv.org/abs/2412.00631}
\bibitem{ref161} INSTRUCTIR: A Benchmark for Instruction Following of Information Retrieval Models. arXiv:2402.14334, 2024. \url{https://arxiv.org/abs/2402.14334}
\bibitem{ref162} Data-Efficiency with a Single GPU: An Exploration of Transfer Methods for Small Language Models. arXiv:2210.03871, 2022. \url{https://arxiv.org/abs/2210.03871}
\bibitem{ref163} RecRanker: Instruction Tuning Large Language Model as Ranker for Top-k Recommendation. arXiv:2312.16018, 2023. \url{https://arxiv.org/abs/2312.16018}
\bibitem{ref164} The Poison of Alignment. arXiv:2308.13449, 2023. \url{https://arxiv.org/abs/2308.13449}
\bibitem{ref165} INSTRUCTEVAL: Towards Holistic Evaluation of Instruction-Tuned Large Language Models. arXiv:2306.04757, 2023. \url{https://arxiv.org/abs/2306.04757}
\bibitem{ref166} Is It Good Data for Multilingual Instruction Tuning or Just Bad Multilingual Evaluation for Large Language Models?. arXiv:2406.12822, 2024. \url{https://arxiv.org/abs/2406.12822}
\end{thebibliography}

\end{document}
